\subsection{Modelling physics in $PGA$}
\label{subsection:modelling_physics}

The key idea is to represent physical entities using objects represented in the underlying
algebra. They have a geometric interpretation and can be manipulated by applying algebraic
operations. This is possible since the algebraic objects we use in geometric algebra,
especially the blades resulting from outer products between vectors, represent subspaces
of the overall space and thus allow for direct algebraic manipulation. At the same time,
they represent or can be attributed to objects of geometric or physical meaning. This
enables implementing geometrically and physically meaningful manipulations by applying
algebraic operations on those subspaces or between them. \\

Euclidean geometric algebra ($EGA$) is limited to objects that contain the origin. Thus,
it can be used to directly model vectors (i.e. directions), lines, and planes through the
origin. In addition, it can model complex numbers, quaternions and thus rotations in $2D$
and $3D$ well. To realize them, it uses sandwich products based on the geometric product
for calculating the transformation resulting from two consecutive reflections either
across lines ($2D$ case) or across planes ($3D$ case). The resulting transformations are
rotations. The advantage of using geometric algebra to model those rotations is the
intuitive modelling, due to geometric meaning of the objects and the fact that the
sandwich products can be used to transform different objects, like vectors, bivectors or
rotation operators using the same sandwich product formulas in a unified fashion
regardless of the type of object to be transformed. Despite those limitations $EGA$ is
already a very powerful tool to express physical ideas as shown by David Hestenes in
\cite{David_Hestenes_NFCM:1986}. In this work he reformulates classical mechanics using
geometric algebra mainly with $EGA$ as a tool. But he is also showing that even more is
possible, e.g. for formulating the motion of rigid bodies, clearly pointing to $PGA$ in
his chapter 7 on rigid body mechanics, even if he does not call that $PGA$ yet. This is
already looking extremely similar to what is shown later in this section. His book is a
rich source of applications of geometric algebra for many application areas of
mechanics.\\

Projective geometric algebra ($PGA$) is indeed capable to model objects that do not
necessarily contain the origin. Thus, it can be used to directly model points, directions,
lines, and planes that may or may not contain the origin. In contrast to $EGA$,
translations can additionally be implemented in $PGA$ directly as orthogonal
transformation using a sandwich product. The operation is realized by two consecutive
reflections across parallel lines ($2D$) or planes ($3D$). It is based on the regressive
form of the geometric product. The reason for this deviation from $EGA$, which uses the
geometric product for that purpose, is that the Euclidean isometries (rotation,
translation, and screw motion) can only be performed in a space that is dual to the space
we use for modelling our projective objects (for a more in depth explanation
see~\cite{Lengyel_pga-illuminated:2024}). The use of the regressive geometric product on
these objects is equivalent to first transforming the object into dual space by applying
the complement operation, then using the geometric product in that space for the intended
transformation and finally transforming the resulting transformed object back into the
original space, by applying the inverse complement operation. Thus, the sandwich product
transformation formulas in $PGA$ only differ from the formulas used in $EGA$ by using the
regressive geometric product instead of the geometric product, but are otherwise the same.
\\

The additional degree of freedom in $PGA$ is possible due to embedding the Euclidean space
of dimension $n$ into a projective space of dimension $n+1$, e.g. two-dimensional
Euclidean space is modelled by embedding it in a three-dimensional projective space, and
three-dimensional Euclidean space is modelled in a four-dimensional projective space. In
the $n+1$-dimensional projective space all algebraic objects contain the origin (exactly
as in $EGA$), but in the modelled $n$-dimensional Euclidean space the objects do not need
to contain the origin any longer. They can be located anywhere in the modelled space
including positions at infinity. This can be realized in projective geometry, since
objects of dimension $n+1$ are projected into the modelled space resulting in objects of
dimension $n$ after the projection. A projective point (which models a scalar or
$0D$-object) is a vector representing a ray through the origin in the embedding space (a
$1D$-object), and a $1D$-object, like a line in modelled space, is created by a bivector
in the embedding projective space (a $2D$-object), and so on for higher-grade objects. \\

Due to its larger expressiveness $PGA$ is suitable to model physics of mass points and
rigid bodies including all isometries. For example, it can model position, velocity,
acceleration, force, torque, momentum, angular momentum, etc., and thus it can be used to
describe kinematics as well as dynamics of mass points and rigid bodies. The approach used
here is inspired by \cite{Plane-based_PGA_Dorst-DeKennik:2022} and
\cite{Dynamics_in_PGA_plane-based_Dorst-DeKennik:2023}, but it does not use the
plane-based ansatz with planes as vectors, which are based on dual representations of the
objects to be modelled. It rather uses the approach as described by
\cite{Lengyel_pga-illuminated:2024} using projective geometric algebra based on direct
representations of points, lines and planes. Lengyel also gives reasons for choosing the
latter approach and explains the background in more detail compared to the short
indications given here. \\

Objects modelled in a projective space can be multiplied by a scalar value without
changing the \emph{geometric} meaning of the object in the modelled Euclidean space, i.e.
its interpretation after the projection. Equality of modelled geometric objects is thus
often only defined up to a scalar factor. This kind of equality up to a scalar factor is
defined as congruence (see e.g. \cite{Browne_Grassmann-Algebra_Vol1:2012}). Therefore, it
is important for modelling of physical facts that unitization is used consciously to avoid
false interpretations. For example, a heavy point can be modelled by multiplying a
\emph{unitized} point with the scalar value of its mass, but of course the multiplication
with $m$ must happen after unitization in order to keep the scalar factor $m$ in the
calculation. Other geometric objects need to be unitized in order to focus on the intended
geometric operation and avoid additional scaling, i.e. the rotation operators. They must
be unitized before use as well. Thus, it is advisable to think through the geometric
operations and their physical interpretation well to avoid surprises and to be consistent.


\newpage

\subsubsection{Modelling force and torque}
\label{subsubsection:modelling_force_and_torque}

\cite{Browne_Grassmann-Algebra_Vol1:2012} (see sections 4.1-4.5) is a very good source
showing how to use Grassmann algebra and projective geometry in order to model forces and
moments. Classical approaches based on vector algebra only model forces by providing
vectors of adequate magnitude and direction to represent the force. However, the
physically very meaningful ``line of action'', which determines which torque is created by
the force, is not modelled mathematically in the classical approach at all. It rather
needs to be derived from the description of the context in order to know where the force
is actually acting. Thus, force and torque are not directly interlinked in the classical
description, but must be modelled as separate physical effects.\\

This changes in $PGA$. Any force can be directly modelled by its line of action, which
encodes the force \emph{and} the torque generated by the force at the origin within one
algebraic object. A force $f$ is modelled by a line $F$. If $f$ is the force vector of
corresponding magnitude and direction, and $Q$ is any unitized projective point on its
line of action then $F$ is an algebraic object representing the force $f$ as well as the
resulting line of action directly:
\begin{equation}
F = Q \wedge f
\label{eq:line_of_action}
\end{equation}
The resulting bivector $F$ represents a line, a geometric object which contains all
physically relevant information of the force $f$ and the torque generated by $f$. For that
reason~\cite{Dynamics_in_PGA_plane-based_Dorst-DeKennik:2023} call this object ``forque''.
It can directly be used to recover the force vector and to calculate the resulting free
torque caused by the force either at the origin or at any point in space. \\

But first, let's have a closer look how this is accomplished. We can separate $F$ into its
components. Each unitized projective point $Q$ is the sum of the origin $O$ and a vector
$q$ parallel to the projective plane (in $2D$) or projective space (in $3D$) pointing from
$O$ to $Q$, such that $Q = O + q$ (``projective point split''). Thus, $F$ can be split into
two bivector components
\begin{equation}
    F = Q \wedge f = (O + q) \wedge f =
    \underbrace{O \wedge f}_{\text{=force}} + \underbrace{q \wedge f}_{=torque}
    = \weight{F} + \bulk{F}
\label{eq:line_of_action_components}
\end{equation}
where the first term encodes the force within the weight components of the bivector, while
the second term encodes the resulting free torque relative to the origin caused by the
action of $f$ at $Q$ -- or alternatively anywhere along the line of action -- within the
bulk components of the bivector. \\

To see this even more clearly, let's look at the component equations of both terms written
as bivectors for the case of $2D$ $PGA$ (the corresponding basis bivectors are shown on
the right-hand side in square brackets for reference):
\begin{subeqnarray}
    O \wedge f & = &
    \begin{pmatrix}
        0 \\ 0 \\ 1 
    \end{pmatrix}
    \wedge
    \begin{pmatrix}
         f_x \\ f_y \\ 0 
    \end{pmatrix}
    =
    \begin{pmatrix}
         f_x \\ f_y \\ 0 
    \end{pmatrix}
    \qquad \qquad \quad \;\;\;\;
    \begin{bmatrix}
         \bv{31} \\ \bv{32} \\ \bv{12} 
    \end{bmatrix}
    \slabel{eq:pga2dp_force_line_force_component} \\[10 pt]
        q \wedge f & = &
    \begin{pmatrix}
        q_x \\ q_y \\ 0 
    \end{pmatrix}
    \wedge
    \begin{pmatrix}
         f_x \\ f_y \\ 0 
    \end{pmatrix}
    =
    \begin{pmatrix}
         0 \\ 0 \\ q_x f_y - q_y f_x
    \end{pmatrix}
    \qquad
    \begin{bmatrix}
         \bv{31} \\ \bv{32} \\ \bv{12} 
    \end{bmatrix}
    \slabel{eq:pga2dp_force_line_torque_component}
\end{subeqnarray}
where equation (\ref{eq:pga2dp_force_line_force_component}) shows how the force is encoded
in the first two components of the bivector and  equation
(\ref{eq:pga2dp_force_line_torque_component}) shows how the torque is encoded in the third
component of the bivector $F$. \\

For $3D$ $PGA$ the encoding is similar, but now requires six bivector components instead
of three (the corresponding basis bivectors are shown on the right-hand side in square
brackets for reference again):
\begin{subeqnarray}
    O \wedge f & = &
    \begin{pmatrix}
        0 \\ 0 \\ 0 \\ 1 
    \end{pmatrix}
    \wedge
    \begin{pmatrix}
         f_x \\ f_y \\ f_z \\ 0 
    \end{pmatrix}
    =
    \begin{pmatrix}
         f_x \\ f_y \\ f_z \\ 0 \\ 0 \\ 0
    \end{pmatrix}
    \qquad \qquad \quad \quad \, %\;\;
    \begin{bmatrix}
         \bv{41} \\ \bv{42} \\ \bv{43} \\ \bv{23} \\ \bv{31} \\ \bv{12} 
    \end{bmatrix}
    \slabel{eq:pga3dp_force_line_force_component} \\[10 pt]
        q \wedge f & = &
    \begin{pmatrix}
        q_x \\ q_y \\ q_z \\ 0 
    \end{pmatrix}
    \wedge
    \begin{pmatrix}
         f_x \\ f_y \\ f_z \\ 0 
    \end{pmatrix}
    =
    \begin{pmatrix}
         0 \\ 0 \\ 0 \\ q_y f_z - q_z f_y \\ q_z f_x - q_x f_z \\ q_x f_y - q_y f_x
    \end{pmatrix}
    \qquad
    \begin{bmatrix}
         \bv{41} \\ \bv{42} \\ \bv{43} \\ \bv{23} \\ \bv{31} \\ \bv{12} 
    \end{bmatrix}
    \slabel{eq:pga3dp_force_line_torque_component}
\end{subeqnarray}

The first three bivector components model the force while the last three components model
the torque. The force vector $f$ itself can be retrieved from $F$ by using the attitude
function $att()$:
\begin{equation}
    f = att(F) = F \vee \overline{\bv{n}} = att(\weight{F})
                                          = \weight{F} \vee \overline{\bv{n}}
    \label{eq:force_attitude}
\end{equation}
with $\bv{n}$ denoting the basis vector of the projective dimension that squares to zero
in the respective algebra (\bv{3} for $2D$ and \bv{4} for $3D$). The torque generated
relative to the origin $O$ is encoded in the bulk of the bivector:
\begin{equation}
    M_{O} = bulk(F) = \bulk{F}
    \label{eq:force_torque_origin}
\end{equation}
while the resulting torque from the application of $f$ at any given point $R$ in space can
be retrieved from $F$, either using any known point $Q$ on $F$\footnote{Which is always
available from the function $sup(F)$. It provides the support point on $F$ with the
shortest distance to the origin $O$ by orthogonally projecting the origin onto $F$.}, or
alternatively by projecting $R$ onto $F$ and using the perpendicular distance vector
$d_\perp$ from $R$ towards the line $F$ in order to calculate the resulting free torque
$M_F(R)$ at the position $R$ as
\begin{subeqnarray}
    M_F(R) & = &  \bulk{\left[ (Q-R) \wedge f \right]} = \bulk{\left[(Q-R) \wedge att(F)
             \right]} \slabel{eq:action_line_moment_simple} \\
    M_F(R) & = & \bulk{\left[ d_\perp \wedge f \right]} = \bulk{\left[ (F \vee (R \wedge
             \rightweightdual{F}) - R) \wedge att(F) \right]}
    \slabel{eq:action_line_moment_projection}
\end{subeqnarray}
where equation~(\ref{eq:action_line_moment_simple}) is the formula with the least
computational effort using a known point $Q$ on $F$, while
equation~(\ref{eq:action_line_moment_projection}) uses the orthogonal projection of $R$
onto $F$ to calculate the perpendicular distance vector $d_\perp$. \\

The resulting total force and torque acting on an object are given by:
\begin{equation}
    F_{total} = \sum_i{F_i}
    \label{eq:total_sum_forque}
\end{equation}
This resulting sum will be used later for connecting kinematics and dynamics. \\

Things to remember:
\begin{itemize}

    \item Adding forces represented as projective lines creates resulting force lines
    located at the right position to get equivalent moments automatically. When two
    parallel forces are added the resulting force line will be located between such that
    it creates a torque equivalent to the torque initially generated by both forces. It
    creates an equivalent force system reducing multiple forces and moments acting on a
    rigid body to a single resultant force and a resultant moment about a specific point
    while maintaining the same external effect on the rigid body.

    \item Thus, the sum of forces and their induced moments can be directly created by
    the sum of their respective lines of action, i.e. by summing the bivectors
    representing those lines. The resulting lines of action contain all relevant
    information incl. the resulting vector sum of the forces and the resulting moment
    values from the summation. There is no need to calculate these separately. The result
    will always be a resulting bivector that potentially contains a bound vector and a
    resulting moment which creates a torque.
    (see~\cite{Browne_Grassmann-Algebra_Vol1:2012}, p. 176).

    \item Neither a vector nor a bivector are located anywhere. They can be moved around
    freely. A line however represents a bound vector that is bound by its line. Thus, the
    vector can only be moved along its line of action without changing the geometric and
    physical meaning (see~\cite{Browne_Grassmann-Algebra_Vol1:2012}, p. 169).

    \item A free torque is modelled by a bivector with force components being zero, i.e.
    the resulting force vector (e.g. from summing several forces acting on an object) has
    no direction and zero length, but the bivector still might contain a resulting torque
    value.

    \item The torque calculated from either equation (\ref{eq:action_line_moment_simple})
    or (\ref{eq:action_line_moment_projection}) is necessarily a free torque as can be
    seen from the formulas. The origin does not occur in the result anymore due to the
    difference of the two projective points involved.

\end{itemize}

% The geometrical situation is depicted in figure XXX for a two-dimensional case.


\newpage

\subsubsection{Modelling linear and angular momentum}
\label{subsubsection:modelling_momentum}

Momentum in the classical approach is split into linear momentum $\vec{p} = m\vec{v}$ and
angular momentum $\vec{L} = \vec{x} \times \vec{p} = \vec{x} \times m \vec{v}$. Both are
represented as vectors, but they represent entities with different physical meaning. So
they cannot be combined into one object in the classical approach, e.g. by adding them. \\

The $PGA$ approach to momentum is to define both, the linear and the angular aspects,
within one algebraic entity. This is in complete analogy to what we have seen in the last
subsection for force and moment resulting in its line of action modelled by a bivector.

\paragraph{Single point mass:}
The momentum $P$ of a point mass $m$ at a projective (unitized) position $X$ that moves
with velocity $\dot{X}$ is defined as a bivector $P$ as
\begin{equation}
    P = X \wedge m \dot{X}
    \label{momentum_definition}
\end{equation}
It defines the overall momentum at position $X$ including linear and angular momentum in
the bivector $P(X)$, which can be shown by splitting the projective point $X$ into its
components $X = O + x$ (``projective point split''). In the $3D$-case with the origin $O =
O_{3dp} = (0,0,0,1)$ and the direction vector $x$, where $x = (X_x, X_y, X_z, 0)$, i.e.
specifically with $X_w = 0$, which represents a direction towards a point at infinity.
This results in
\begin{subeqnarray}
    P(X) &=& X \wedge m \dot{X} = (O + x) \wedge m \dot{X} \nonumber \\
         &=& O \wedge m \dot{X} + x \wedge m \dot{X} \nonumber \\
         &=& O \wedge m \dot{(O + x)} + x \wedge m \dot{(O + x)} \nonumber \\
         &=& O \wedge m \dot{x} + x \wedge m \dot{x} \\
         &=& \underbrace{\quad\weight{P}\quad}_{\text{linear}}
             \;+\underbrace{\quad\bulk{P}\quad}_{\text{angular}}
\end{subeqnarray}
$P$ defines an instantaneous momentum line where the weight part $\weight{P}$ of the
bivector $P$ encodes the linear momentum and the bulk part $\bulk{P}$ encodes the angular
momentum with respect to the origin.\footnote{The ``projective point split'' works in this
way, since the origin is a distinguishable point. This is in contrast to $EGA$ or the
classical approach where the origin is a null-vector, which has no lasting effect if it is
added to another point. This is one reason why $PGA$ can combine linear and angular
aspects as separate physical effects in a single algebraic object. The other is that the
required 6 degrees of freedom can be encoded in a bivector.}

\paragraph{Recovering linear and angular momentum:}
The physical aspects encoded in the momentum line $P$ can be retrieved in complete analogy
to the case of force and torque in subsection
\ref{subsubsection:modelling_force_and_torque}: \\

The linear momentum vector $p$ itself can be retrieved from $P$ by using the attitude
function $att()$:
\begin{equation}
    p = m v = m \dot{x} = att(P) = P \vee \overline{\bv{n}}
                        = att(\weight{P}) = \weight{P} \vee \overline{\bv{n}}
    \label{eq:linear_momentum_from_attitude}
\end{equation}
with $\bv{n}$ denoting the basis vector of the projective dimension that squares to zero
in the respective algebra (\bv{3} for $2D$ and \bv{4} for $3D$). \\

The angular momentum $L$ generated relative to the origin $O$ is encoded in the bulk of
the bivector:
\begin{equation}
    L_{O} = x \wedge m \dot{x} = bulk(P) = \bulk{P}
    \label{eq:angular_momentum_wrt_origin}
\end{equation}
while the resulting angular momentum relative to any given point $R$ in space can be
retrieved from the momentum line $P$ by
\begin{equation}
    L(R) = \bulk{\left[ (X-R) \wedge m\dot{X} \right]}
         = \bulk{\left[ (X-R) \wedge att(P) \right]}
           \slabel{eq:action_momentum_line_simple}
\end{equation}
using $X$ as a known point on $P$. This provides the resulting angular momentum $L(R)$
with respect to $R$. The momentum line encodes the resulting momentum including its linear
and angular aspects relative to any arbitrarily chosen point. \\

The equations (\ref{eq:linear_momentum_from_attitude}) and
(\ref{eq:angular_momentum_wrt_origin}) given above permit to separate the linear and
angular effects as in the classical approach. However, these aspects need not be separated
at all and can be calculated together in $PGA$.\\

Using eqn.~(\ref{momentum_definition}) and the definition $X=O + x$ with $X$ as unitized
projective point, $O$ as the chosen origin and $x$ as the direction vector pointing from
$O$ to $X$ we can write
\begin{equation}
    P_O(X) = X \wedge m \dot{X} = (O + x) \wedge m \dot{(O + x)}
    \label{momentum_written_with_O}
\end{equation}
Compared to eqn.~(\ref{momentum_definition}) an index $O$ was added to $P$, since the
formula (\ref{momentum_written_with_O}) seems to confirm that the momentum depends on the
choice of the origin. Now it is important to demonstrate  that this is \emph{not} the
case. We are going to demonstrate that the described physics does not depend on the chosen
origin and not on the resulting coordinates used for its description. \\

We use a transformation of the standpoint from which we describe the
physical situation without actually moving the coordinate system (we still want to use the
chosen coordinate system to describe the vectors and bivectors used): We go to a virtual
origin $O' = Y$, but describe the situation in the same coordinate system. This results in
describing the origin as $O = Y - y$ as seen from $Y$. Inserting that definition into
eqn.~(\ref{momentum_written_with_O}) and using the geometric vector identity
\begin{equation}
X = Y + (X-Y) = Y + (O+x-(O+y)) = Y + (x-y)
\label{eq:identity_transformation_X_Y}
\end{equation}
we express $P(X)$ as viewed from $Y$ by reformulating the equation for $P$ relative to the
origin $O$
\begin{eqnarray}
    P_O(X) &=& X \wedge m \dot{X} = (O + x) \wedge m \dot{(O + x)} \nonumber \\
           &=& (Y - y + x) \wedge m \dot{(Y - y + x)} \nonumber \\
           &=& (Y + (x-y)) \wedge m \dot{(Y + (x - y))} = P_Y(X)
\end{eqnarray}
where the last equation obviously is the same as eqn.~(\ref{momentum_written_with_O}), but
here Y takes the role of $O$ and $(x-y)$ takes the role of the vector $x$ in the old
description. The dot used for the time-derivative refers to the whole expression in
parentheses. So the same physical situation was described from another viewpoint. However,
the algebraic object $P$ is unchanged, i.e. $P_O(X) = P_Y(X)$. This clearly shows that $P$
does not depend on the chosen origin and that eqn.~(\ref{momentum_written_with_O}) is
indeed a coordinate free definition.

\paragraph{Mass distributions:}
The total momentum of a set of mass points $m_i X_i$ is defined as the sum of all momentum
blades.
\begin{equation}
    P_{total} = \sum_i{P_i}
    \label{eq:total_sum_momentum}
\end{equation}
$P_{total}$ will always be a momentum blade in $2D$ $PGA$, while in $3D$ $PGA$ (a
four-dimensional algebra) it could also result in a general bivector that might not be
factorizable into a blade (see \cite{Browne_Grassmann-Algebra_Vol1:2012}). That would be
the case, if $P \wedge P \ne 0 \; \ps$, i.e. when the result is a pseudoscalar that is not
equal to a zero multiple of the pseudoscalar \bv{1234}. This test can always be used to
determine whether a bivector is a blade in $3D$ $PGA$. \\

Using discrete points $X_{w,i}$ of the world frame indexed by $i$ that can be moved by a
motor $M(t) = exp_\veedot[\frac{1}{2} B_w(t)] \veedot M_0$ in the world frame as by
eqn.~(\ref{eq:pga_sandwich_tranformation}) $X_{w,i} = M \veedot X_{0,i} \veedot \rrev{M}$
results in a local speed of $\dot{X}_{w,i} = \mathrm{rcmt}(\Omega_w,X_{w,i})$ as shown in
eqn.~(\ref{eq:rcmt_pga_world}). Therefore, we can write the total momentum $P_{w}$ in the
world frame as
\begin{eqnarray}
    P_{w} &=& \sum_i{P_i} \nonumber \\
          &=& \sum_i X_i \wedge m_i \dot{X}_i \nonumber \\
          &=& \sum_i X_i \wedge m_i \; \mathrm{rcmt}(\Omega_w,X_i) \nonumber \\
    P_{w} &\equiv& I_w[\Omega_w]
                \label{eq:inertia_map_world_frame}
\end{eqnarray}
where eqn.~(\ref{eq:inertia_map_world_frame}) is a linear function of $\Omega_w$ that
maps the bivector rate in the world frame $\Omega_w$ to the total momentum $P_w$. It is a
function that takes a bivector as an argument and returns a bivector (the ``inertia map''
$I_w$). It describes the influence of the mass distribution of the points on the movement
via the combined Newton/Euler law
\begin{equation}
    \dot{P}_{w}=F_{w}
    \label{eq:Newton_Euler_law}
\end{equation}
(for the total force in the \emph{world frame} $F_{total} = F_{w}$ see
eqn.~(\ref{eq:total_sum_forque})). If we go from discrete mass points to a continuous mass
distribution, the sum has to be replaced by an integral over $dm = \rho dV$ (see following
subsections on inertia maps for more details). \\

The equation $\dot{P}_{w} = X_w \wedge m \ddot{X}_w = F_{w}$ cannot be integrated directly
since the second derivative $\ddot{X}_w$ cannot be isolated on the left-hand side of the
equation as required for direct integration. To circumvent this we introduce another
variable that can be isolated, enabling us to solve the ODE. For this purpose we use the
rate of change $\Omega$ (a vector in $2D$ and a bivector in $3D$). The integration of that
variable will then allow for computation of position, speed, etc. as a separate step.

\paragraph{Rate of change:}
Calculating the derivative $\dot{P}_w$ in the world frame in general yields two terms
\begin{equation}
    \dot{P}_w = \dot{I}_w[\Omega_w] + I_w[\dot{\Omega}_w]
    \label{eq:momentum_derivative_world_frame}
\end{equation}
since in the world frame the inertia $I_w[\Omega_w]$ map will change over time due to
moving points as well as due to the rate of change $\Omega_w$. Therefore, it is reasonable
to calculate the inertia map in the body frame, since it is constant in that frame, if the
body frame is chosen to be fixed to the modelled rigid object. This avoids the effort to
additionally calculate the time derivative of the inertia map (the first term on the
right-hand side above).

\paragraph{Body frame:}
If we want to express the physical problem in the body frame, we have to transform the
components of eqn.~(\ref{eq:inertia_map_world_frame}) accordingly. Using
eqn.~(\ref{eq:world_to_body_pga}), i.e. $I_b = \rrev{M} \veedot I_w \veedot M$ and
eqn.~(\ref{eq:body_to_world_pga}) $\Omega_w  = M \veedot \Omega_b \veedot \rrev{M}$ then
results in
\begin{subeqnarray*}
    P_{w} &=& I_w[\Omega_w] \\
    \underbrace{\rrev{M} \veedot P_{w} \veedot M}_{P_{b} } &=&
    \underbrace{\rrev{M} \veedot I_w[\rrev{M} \veedot \Omega_w \veedot M] \veedot M}_
    {I_b[\Omega_b]}
\end{subeqnarray*}
which finally results in an expression for the inertia map
\begin{equation}
    P_b = I_b[\Omega_b] \equiv \rrev{M} \veedot I_w[ M \veedot \Omega_b \veedot \rrev{M} ]
    \veedot M
    \label{eq:inertia_map_body_frame}
\end{equation}

The same transformation must be applied to $F_w$ for the right-hand side of the
differential equation $\dot{P}_w = F_w$. Therefore, we get in the body frame $F_b =
\rrev{M} \veedot F_w \veedot M$, resulting in $\dot{P}_b = F_b$. \\

The time derivative in the body frame is simpler, as long as we restrict ourselves to
rigid bodies, resulting in only one term, since the rate of change $\Omega_b$ is the only
time dependent function in the body frame
\begin{equation}
    \dot{P}_b = I_b[\dot{\Omega}_b]
    \label{eq:momentum_derivative_body_frame}
\end{equation}
otherwise $\dot{I}_b[\Omega_b]$ would have to be covered, as in
eqn.~(\ref{eq:momentum_derivative_world_frame}) in the world frame. \\

The mapping to the classical inertia can be done when splitting $\Omega_b$ relative to
the centroid of the body with a projective point split on $\Omega_b$ and on
$I_b[\Omega_b]$ as shown by \cite{Dynamics_in_PGA_plane-based_Dorst-DeKennik:2023} on page
25f and adjusted to our approach using regressive geometric products.

\paragraph{Reference points:}
Now we introduce some typically used reference points to describe kinematics via the
bivector rate of change $\Omega$ and dynamics via the total momentum $P$:
\begin{itemize}
    \item \emph{Fixed body reference point} $O_b$ used as origin of the body frame. When
    the body is moved by the motor $M$ it will move in the world frame to $O_{b,w} = M
    \veedot O_b \veedot \rrev{M}$.
    \item \emph{Relative locations $r_i$ versus $O_b$} are positions in the body frame.
    Their projective counterpart in the body frame is $R_i = O_b + r_i$. When transformed
    to the world frame we get $X_i$ in the world frame by the usual transformation $X_i =
    M \veedot R_i \veedot \rrev{M} = M \veedot (O_b + r_i) \veedot \rrev{M}$.
    \item \emph{Center of mass} in the body frame is chosen to be located at $O_b$, such
    that $\sum_i{m_i r_i} = \vec{0}$. Thus, the origin of the body frame is located at the
    center of mass. The center of mass in the world frame $C_w$ then is $C_w \equiv
    \sum_i{m_i X_i} / \sum_i{m_i} = \sum_i{m_i (M \veedot (O_b + r_i) \veedot \rrev{M})} /
    \sum_i{m_i} = M \veedot O_b \veedot \rrev{M}$.
\end{itemize}

The last missing piece for dynamics is the time derivative of $P_{w}$ expressed as time
rate of change $\Omega_b$ of the inertia in the body system. Applying
eqn.~(\ref{eq:rcmt_pga_world}) in the body frame, which is valid for any rigid object
moved by a motor $M$, we get
\begin{equation}
    \frac{d}{dt} P_b = \frac{d}{dt} I_b[\Omega_b] =
    \mathrm{rcmt}(\Omega_b,I_b[\Omega_b]) + I_b[\dot{\Omega}_b]
\end{equation}
where the last term arises since $\Omega_b$ is also time-dependent. Solving
eqn.~(\ref{eq:Newton_Euler_law}), but in the body frame, we get
\begin{equation}
    \dot{P}_b = \mathrm{rcmt}(\Omega_b,I_b[\Omega_b]) + I_b[\dot{\Omega}_b] = F_b
\end{equation}
resolving this for $\dot{\Omega}_b$ we get
\begin{eqnarray}
    I_b[\dot{\Omega}_b] &=& F_{b} - \mathrm{rcmt}(\Omega_b,I_b[\Omega_b]) \\[5pt]
    \dot{\Omega}_b &=& I_b^{-1}\left[ F_{b} - \mathrm{rcmt}(\Omega_b,I_b[\Omega_b])
    \right]
\end{eqnarray}

Together with the ansatz in the body system eqn.~(\ref{eq:motor_pga_body_frame}) that can
be used to calculate M(t) when we know $\Omega_b$ at a given time $t$ from
\begin{equation*}
   M(t) = M_0 \veedot \exp_\veedot(\frac{1}{2} B_b(t))
        = M_0 \veedot \exp_\veedot(\frac{1}{2} \Omega_b t) 
\end{equation*}
and the rate of change in the body frame according to
eqn.~(\ref{eq:motor_rate_body_frame_pga}) we get
\begin{equation}
    \dot{\Omega}_b = I_b^{-1}\left[ F_{b} - \mathrm{rcmt}(\Omega_b,I_b[\Omega_b]) \right] 
    \label{eq:ode_system_body_frame}
\end{equation}
as an ODE that can be solved as initial value problem. Due to $\dot{\Omega}_b =
\ddot{B}_b$ it actually is a $2^{nd}$ order ODE (like the classical Newton/Euler equations
that can be integrated directly from the instantaneous acceleration and the angular
acceleration, respectively, in order to calculate speed and position). Here we rather
integrate from $\ddot{B}_b$ in order to calculate $\dot{B_b} = \Omega_b$ and $B_b$. Thus,
we require initial values for $\Omega_b$ and $B_b$. The actual positions of the points of
the rigid body can then be computed from eqn.~(\ref{eq:pga_sandwich_tranformation})
$X_w(t) = M(t) \veedot X_0 \veedot \rrev{M}(t)$. The speed of the points can be calculated
either from eqn.~(\ref{eq:rcmt_pga_world}) $\dot{X}_w = \mathrm{rcmt}(\Omega_w,X_w)$ or
eqn.~(\ref{eq:rcmt_pga_world_incl_X0_of_t}) $\dot{X}_w = \mathrm{rcmt}(\Omega_w,X_w)+
(\dot{X_b})_w$ -- which one is applicable depends on whether the rigid body is fixed in
the body frame or moves in the body frame. The expressions in the world frame can of
course be expressed in the body frame as well using eqn.~(\ref{eq:body_to_world_pga})
$\Omega_w = M \veedot \Omega_b \veedot \rrev{M}$ and $X_w = M \veedot X_b \veedot
\rrev{M}$.

\paragraph{Inertia map ($2D$):}
For that purpose we need the inertia map $I_b[\Omega_b]$, and it's inverse to set up the
right-hand side of the ODE system given in equations~(\ref{eq:ode_system_body_frame}). We
start with the $2D$ case where the rate of change $\Omega_b$ is a vector, as shown at the
end of subsection \ref{subsection_modelling_motion}. We assume discrete points $X_i$ and
use the wedge product $X \wedge Y$ and the regressive commutator product term $m \;
\mathrm{rcmt}(\Omega_b,X)$ to express the total momentum $P$ written in component form:
\begin{subeqnarray}
    X \wedge Y & = &
    \begin{pmatrix}
        X_x \\ X_y \\ X_z  
    \end{pmatrix}
    \wedge
    \begin{pmatrix}
         Y_x \\ Y_y \\ Y_z
    \end{pmatrix}
    =
    \begin{pmatrix}
         X_z Y_x - X_x Y_z \\ X_z Y_y - X_y Y_z \\ X_x Y_y - X_y Y_x
    \end{pmatrix}
    \slabel{eq:wedge_X_Y_2D} \\[10 pt]
    m \; \mathrm{rcmt}(\Omega_b, X) & = & m
    \begin{pmatrix}
         \Omega_y X_z - \Omega_z X_y \\ \Omega_z X_x - \Omega_x X_z \\ 0
    \end{pmatrix} 
\end{subeqnarray}

Now the inertia map for a point $X$ with mass $m$ and rate of change $\Omega$ can be
expressed in the $2D$ case as
\begin{eqnarray}
    I_{2D}[\Omega] &=& X \wedge m \; \mathrm{rcmt}(\Omega,X) \nonumber \\[10pt]
    &=& m 
    \begin{pmatrix}
        X_z (\Omega_y X_z - \Omega_z X_y) \\
        X_z (\Omega_z X_x - \Omega_x X_z) \\
        X_x (\Omega_z X_x - \Omega_x X_z) - X_y (\Omega_y X_z - \Omega_z X_y)
    \end{pmatrix} \nonumber \\[10pt]
    I_{2D}[\Omega] &=& m
    \begin{bmatrix}
        0        &  X_z^2   & -X_y X_z      \\
        -X_z^2   &  0       &  X_x X_z      \\
        -X_x X_z & -X_y X_z & X_x^2 + X_y^2
    \end{bmatrix}
    \begin{bmatrix}
        \Omega_x \\ \Omega_y \\ \Omega_z
    \end{bmatrix}
    \label{eq:inertia_map_components_2D}      
\end{eqnarray}
where eqn.~(\ref{eq:inertia_map_components_2D}) is linear in $\Omega$ as expected. Of
course, the total inertia map needs to be summed for the contributions from each point,
i.e. $I_{total} = \sum_i{I_i}$. Since all points of the rigid body see the same rate of
change the total inertia map can be summed for all mass point contributions independent of
the rate of change and the resulting total inertia map can be inverted either symbolically
or numerically to get its inverse $I_{2D}^{-1}$ as needed to set up the right-hand side of
the system in eqn.~(\ref{eq:ode_system_body_frame}). The inertia map in $2D$ maps the rate
of change vector to the momentum bivector. The inverse inertia map than maps from a
momentum bivector to the rate of change vector.

\paragraph{Inertia map ($3D$):}
Now let's do the same for the $3D$ case, where the rate of change is a bivector. The
component expressions are generated from the corresponding product tables in
\path{ga_prdxpr} and integrated into the products in the ga library. 
\begin{subeqnarray}
    X \wedge Y & = &
    \begin{pmatrix}
        X_x \\ X_y \\ X_z  \\ X_w
    \end{pmatrix}
    \wedge
    \begin{pmatrix}
         Y_x \\ Y_y \\ Y_z \\ Y_w
    \end{pmatrix}
    =
    \begin{pmatrix}
         X_w Y_x - X_x Y_w \\
         X_w Y_y - X_y Y_w \\
         X_w Y_z - X_z Y_w \\
         X_y Y_z - X_z Y_y \\
         X_z Y_x - X_x Y_z \\
         X_x Y_y - X_y Y_x
    \end{pmatrix}
    \slabel{eq:wedge_X_Y_3D} \\[10 pt]
    m \; \mathrm{rcmt}(\Omega_b,X) & = & m
    \begin{pmatrix}
         \Omega_{vy} X_z - \Omega_{vz} X_y + \Omega_{mx} X_w \\
         \Omega_{vz} X_x - \Omega_{vx} X_z + \Omega_{my} X_w \\
         \Omega_{vx} X_y - \Omega_{vy} X_x + \Omega_{mz} X_w \\
         0
    \end{pmatrix} 
\end{subeqnarray}
 
Now the inertia map for a point $X$ with mass $m$ and rate of change $\Omega$ can be
expressed in the $3D$ case as
\begin{align}
    &I_{3D}[\Omega] = X \wedge m \; \mathrm{rcmt}(\Omega,X) = \nonumber \\[10pt]
    &m 
    \begin{pmatrix}
        X_w (\Omega_{vy} X_z - \Omega_{vz} X_y + \Omega_{mx} X_w) \\
        X_w (\Omega_{vz} X_x - \Omega_{vx} X_z + \Omega_{my} X_w) \\
        X_w (\Omega_{vx} X_y - \Omega_{vy} X_x + \Omega_{mz} X_w) \\
        X_y (\Omega_{vx} X_y - \Omega_{vy} X_x + \Omega_{mz} X_w) 
      - X_z (\Omega_{vz} X_x - \Omega_{vx} X_z + \Omega_{my} X_w) \\
        X_z (\Omega_{vy} X_z - \Omega_{vz} X_y + \Omega_{mx} X_w)
      - X_x (\Omega_{vx} X_y - \Omega_{vy} X_x + \Omega_{mz} X_w) \\
        X_x (\Omega_{vz} X_x - \Omega_{vx} X_z + \Omega_{my} X_w) - X_y (\Omega_{vy} X_z -
      \Omega_{vz} X_y + \Omega_{mx} X_w) \end{pmatrix} = \nonumber \\[10pt]
    &m
    \begin{bmatrix}
      0              &  X_z X_w       & -X_y X_w       & X_w^2    &  0       &  0       \\
      -X_z X_w       &  0             &  X_x X_w       & 0        &  X_w^2   &  0       \\
       X_y X_w       & -X_x X_w       &  0             & 0        &  0       &  X_w^2   \\
       X_y^2 + X_z^2 & -X_x X_y       & -X_x X_z       & 0        & -X_z X_w &  X_y X_w \\
      -X_x X_y       &  X_x^2 + X_z^2 & -X_y X_z       &  X_z X_w &  0       & -X_x X_w \\
      -X_x X_z       & -X_y X_z       &  X_x^2 + X_y^2 & -X_y X_w &  X_x X_w &  0
    \end{bmatrix}
    \begin{bmatrix}
        \Omega_{vx} \\ \Omega_{vy} \\ \Omega_{vz} \\
        \Omega_{mx} \\ \Omega_{my} \\ \Omega_{mz}
    \end{bmatrix} 
    \label{eq:inertia_map_components_3D}   
\end{align}
where eqn.~(\ref{eq:inertia_map_components_3D}) is linear in $\Omega$ as expected. Again,
the total inertia map needs to be summed for the contributions from each point, i.e.
$I_{total} = \sum_i{I_i}$. The resulting total inertia map can be inverted numerically to
get its inverse $I_{3D}^{-1}$ as needed to set up the right-hand side of the system in
eqn.~(\ref{eq:ode_system_body_frame}). The inertia map in $3D$ maps the rate of change
bivector to the momentum bivector. The inverse inertia map than maps from a momentum
bivector to the corresponding rate of change bivector.\\

As shown in \cite{Dynamics_in_PGA_plane-based_Dorst-DeKennik:2023} the inertia map is
additive. That means that inertia maps of parts of a body consisting of several parts that
move with the same bivector rate can be added to get the total inertia. \\


\newpage

\subsubsection{Solving the Newton-Euler equations numerically}
\label{subsubsection:solving_newton-euler_numerically}

In classical mechanics the Newton equation states that change of linear momentum $\vec{p}
= m \vec{v}$ of a mass point of mass $m$ moving with speed $\vec{v}$ is driven by an
external force $\vec{f}$ acting on the mass. This results in $\dot{\vec{p}} = m
\dot{\vec{v}} = m \ddot{\vec{x}}_{cm} = \vec{f}$ for translational motion in a system with
constant mass $m$. The center of mass at position $\vec{x}_{cm}$ moves with the velocity
$\vec{v}$ as a result. \\

The analogous equation for rotational motion, attributed to Euler, states that the
temporal change of angular momentum $\vec{L} = \vec{x} \times \vec{p} = \vec{x} \times m
\vec{v}$ is driven by an external angular momentum $\vec{M}$. $\vec{M}$ is given as
$\vec{M} = \vec{r} \times \vec{f}$ with $\vec{r}$ as the vector pointing from the
rotational axis to the point of action of the force $\vec{f}$. The equation for rotational
motion $\dot{\vec{L}} = J \dot{\vec{\omega}} = \vec{M}$ for a rigid body with moment of
inertia $J = \int_{m}{r^2_\perp dm} = \int_{V}{r^2_\perp \rho dV}$ with $r_\perp$ as the
perpendicular distance of the mass point within the rigid body to the momentary rotational
axis. The equation determines the temporal development of the orientation of the rigid
body which freely spins with the angular speed $\nrm{\vec{\omega}}$ around the momentary
rotation axis defined by $\vec{\omega}$  as a result. \\

Due to the same structure the equations of a freely moving body can be written as one
system of equations that can be solved as uncoupled equations (e.g. covering cases with
either translation or rotation exclusively within the identical structure). In the $2D$
case the translational part requires two degrees of freedom (DoF) and the rotational part
needs one DoF resulting in a total of three DoF. It can be written compactly as a system
introducing the inertia matrix $I$ as a linear function that maps the acceleration degrees
of freedom to the corresponding external force and torque components\footnote{Be aware,
that this does not actually couple the equations! We just solve two independent equations
conveniently in parallel to show the principle here. If the movement is actually coupled,
coupling terms must be added to model physics correctly. In the $PGA$ formulation this
coupling is already captured by the $\mathrm{rcmt}$ term in the body-frame equation of
motion~(\ref{eq:ode_system_body_frame}).}:
\begin{equation}
    \underbrace{
    \begin{bmatrix}
        m & 0 & 0 \\
        0 & m & 0 \\
        0 & 0 & J
    \end{bmatrix}
    }_{inertia \; I}
    \begin{pmatrix}
        \ddot{x}_x \\
        \ddot{x}_y \\
        \ddot{\phi} 
    \end{pmatrix}
    =
    \begin{pmatrix}
        {f}_x \\
        {f}_y \\
        {M} 
    \end{pmatrix} 
\end{equation}
Solving such a system requires to isolate the accelerations by multiplying the system from
the left with the inverse of the inertia map $I^{-1}$ resulting in
\begin{equation}
    \begin{pmatrix}
        \ddot{x}_x \\
        \ddot{x}_y \\
        \ddot{\phi} 
    \end{pmatrix}
    =
    \underbrace{
    \begin{bmatrix}
        \frac{1}{m} & 0           & 0 \\
        0           & \frac{1}{m} & 0 \\
        0           & 0           & \frac{1}{J}
    \end{bmatrix}
    }_{inverse \; inertia \; I^{-1}}
    \begin{pmatrix}
        {f}_x \\
        {f}_y \\
        {M} 
    \end{pmatrix} 
    \label{eq:classical_sys_Newton_Euler_2nd_order}
\end{equation}

Both equations result in a system of ordinary differential equation (ODE) of second order
that requires two initial conditions, one for the initial position $\vec{x}(t=0) =
\vec{x}_0$ and the initial speed $\vec{v}(t=0) = \vec{v}_0$ (Newton equation) and one for
the initial orientation $\phi(t=0) = \phi_0$ and the initial angular speed $\omega(t=0) =
\omega_0$ (Euler equation) to be integrated. In the $2D$ case $\phi$ and $\omega$ are
scalar values, while in the $3D$ case they also become vectors each representing three
degrees of freedom.\\

To integrate the equations the system of $2^{nd}$ order must be transferred into a system
of $1^{st}$ order. We have a $2^{nd}$ order ODE for $y(t)$ with $\ddot{y} =
\frac{d^2}{dt^2}(y)$ of the general form
\begin{equation}
    \ddot{\vec{y}} = f(\dot{\vec{y}}, \vec{y}, t)
\end{equation}
The vector $y$ has components that we index with $n$. We use an additional index to
distinguish between the position (index 0) and the speed (index 1) of the respective DoF.
With the substitution
\begin{subeqnarray}
    y_{n,0} &=& y_n \qquad \rightarrow \textit{position} \\
    y_{n,1} &=& \dot{y}_n \qquad \rightarrow \textit{speed}
\end{subeqnarray}
this results in the following system of $1^{st}$ order for each component $n$ of the
vector $\vec{y}$:
\begin{subeqnarray}
    \dot{y}_{n,0} &=& y_{n,1} \\
    \dot{y}_{n,1} &=& f(y_{n,1}, y_{n,0}, t)
\end{subeqnarray}
If we would have to integrate for more than one point, one could easily add an index
at the beginning to extend the system accordingly. However, for now we only focus on
integrating the system at one point (the center of mass) at which we want to know position
and speed as well as angular orientation and angular speed as function of time. The
substitution above transforms the system in
eqn.~(\ref{eq:classical_sys_Newton_Euler_2nd_order}) into a system of $1^{st}$ order as
follows:
\begin{subeqnarray}
    \dot{x}_x &=& v_x \\
    \dot{v}_x &=& \frac{1}{m}f_x \\
    \dot{x}_y &=& v_y \\
    \dot{v}_y &=& \frac{1}{m}f_y \\
    \dot{\phi} &=& \omega \\
    \dot{\omega} &=& \frac{1}{J} M
\end{subeqnarray}
This $1^{st}$ order ODE system $\dot{\vec{u}} = rhs(\vec{u},t)$ can now be solved by
direct integration, using e.g. a Runge-Kutta-Scheme of $4^{th}$ order accuracy for precise
time integration, and using the initial conditions defined above. As a result, we directly
get the positions and velocities as function of time.

\paragraph{Classical body frame vs.\ $PGA$ --- what unifies:} It is instructive to compare
the
single $PGA$ equation of motion~(\ref{eq:ode_system_body_frame}), $\dot{\Omega}_b =
I_b^{-1}\!\left[F_b - \mathrm{rcmt}(\Omega_b, I_b[\Omega_b])\right]$, with the classical
body-frame formulation. Classically one carries two separate vector equations -- Newton's
for $\dot{\vec{v}}$ and Euler's for $\dot{\vec{\omega}}$ (the latter with the nonlinear
gyroscopic term $\vec{\omega}\times\mathbf{I}_b\vec{\omega}$) -- plus a separate
orientation update via a rotation matrix or quaternion. The $PGA$ formulation collapses
all of this into one bivector ODE plus a motor update, as summarised in
Table~\ref{tab:classical_vs_pga}. \\

\begin{table}[!ht]
\centering\small
\begin{tabular}{p{6.0cm}|p{6.5cm}}
\hline
Classical body frame &
$PGA$ equivalent \\
\hline
Two separate equations: Newton + Euler &
One unified equation for $\dot{\Omega}_b$ \\
Coupling term $\vec{\omega}\times\vec{v}$ &
Encoded in $\mathrm{rcmt}(\Omega_b, I_b[\Omega_b])$ \\
Gyroscopic term $\vec{\omega}\times\mathbf{I}_b\vec{\omega}$ &
Also encoded in $\mathrm{rcmt}(\Omega_b, I_b[\Omega_b])$ \\
Inertia tensor $\mathbf{I}_b$ ($3\times3$ or $6\times6$ matrix) &
Inertia map $I_b[\cdot]$ (linear operator on bivectors) \\
Quaternion evolution $\dot{\mathbf{q}} = \frac{1}{2}\mathbf{q}\otimes\vec{\omega}$ &
Motor $M = \exp_\veedot[\frac{1}{2} B_b]$ -- orientation implicit in $B_b$ \\
State: $(\vec{r},\, \mathbf{q},\, \vec{v},\, \vec{\omega})$ &
State: $(B_b,\, \Omega_b)$ -- bivector pair \\
Inertial-frame position from $\dot{\vec{r}} = \mathbf{R}\vec{v}$ &
Inertial-frame points updated via sandwich product
$X(t) = M \veedot X_b(t) \veedot \rrev{M}$ \\
\hline
\end{tabular}
\caption{Classical body-frame vs.\ $PGA$ formulation: explicit correspondence.}
\label{tab:classical_vs_pga}
\end{table}

The key unification is this: where the classical formulation needs one equation for
$\dot{\vec{v}}$ (3 DoF), one for $\dot{\vec{\omega}}$ (3 DoF), and a separate orientation
integration, the $PGA$ ODE handles the first two in a single equation over the bivector
$\dot{\Omega}_b$ (6 DoF) and recovers position and orientation from a motor derived
directly from $\Omega_b$. The single term $\mathrm{rcmt}(\Omega_b, I_b[\Omega_b])$ plays
the combined role of both classical coupling terms ($\vec{\omega}\times\vec{v}$ and
$\vec{\omega}\times\mathbf{I}_b\vec{\omega}$). As described above, the $PGA$ system is
first solved for $B(t)$, $\dot{B} = \Omega(t)$ and $\ddot{B} = \dot{\Omega}(t)$; the
actual points and their velocities then follow in a second step from the motor and the
sandwich transformation. \\

Basic examples for translational and rotational motion are implemented as tests in
\path{ga_test/src/ga_appl_physics_test.hpp} based on the Newton-Euler equations in both
approaches, the classical approach and the approach directly formulated in $PGA$, for very
simple reference cases.

\paragraph{Constrained motion --- rotation about a fixed pivot:}
The body-frame $PGA$ formulation extends naturally to constrained motion. A common case in
$2D$ is a rigid body whose rotation is constrained so that a specific body-frame point
$Q_b$ (the pivot) remains fixed in the world -- for example a plate pendulum fixed at one
of its corners. Two questions are to be answered independently:
\begin{enumerate}
    \item \emph{In which frame are the equations expressed?} (body frame vs.\ world frame)
    \item \emph{About which axis does the body actually rotate?} (center of mass vs.\
    pivot)
\end{enumerate}
The Euler equation $\dot{\Omega}_b = I_{cm}^{-1}[F_b - \mathrm{rcmt}(\Omega_b,
I_{cm}[\Omega_b])]$ describes a \emph{free} rigid body rotating about its own center of
mass. When the pivot is fixed, the center of mass follows a circular arc of radius $r$
around the pivot, with a displacement of the rotation point relative to the center of
mass, effectively increasing the inertia relative to that pivot point. The constraint
reaction also crates centripetal acceleration to keep the center of mass moving along the
arc. This is felt in the body system at the center of mass as additional centrifugal
acceleration. The Steiner (parallel axis) theorem gives the corrected moment of inertia in
$2D$:
\begin{equation}
    I_\text{pivot}[2,2] \;=\; I_{cm}[2,2] + m r^2
    \label{eq:steiner_2d}
\end{equation}
so the angular acceleration about the fixed pivot is $\alpha = \tau_\text{eff} /
I_\text{pivot}[2,2]$, \emph{not} $\tau_\text{eff}/I_{cm}[2,2]$. This is a correction for
the physical rotation axis, not for the choice of frame -- it is required regardless of
whether the equations are expressed in the body frame or the world frame. \\

In the $PGA$ implementation the scalar Steiner correction is built into the inertia
constructor via an optional pivot parameter:
\begin{itemize}
    \item \emph{$2D$ -- pivot is a point} $Q_b$ (grade~1, \texttt{Vec2dp}): \\
    \texttt{get\_plate\_inertia(m, w, h, Q\_b)} updates only $I[2,2]$ by adding
    $m\,r^2$ where $r^2 = Q_{b,x}^2 + Q_{b,y}^2$ is the squared distance from the
    body origin to $Q_b$. All other entries (including off-diagonal coupling terms)
    remain equal to $I_\text{body}$. The default $Q_b = (0,0,1)$ (body origin) gives
    $I_\text{body}$ with no correction.
    \item \emph{$3D$ -- pivot is a line} $L_\text{pivot}$ (grade~2, \texttt{BiVec3dp}): \\
    \texttt{get\_cuboid\_inertia(m, w, h, d, L\_pivot)} computes the foot of the
    perpendicular $P_\text{foot} = (P_x, P_y, P_z)$ from the body origin to
    $L_\text{pivot}$ and updates only the diagonal entries of the lower-left block:
    $I[3,0] \mathrel{+}= m(P_y^2+P_z^2)$, $I[4,1] \mathrel{+}= m(P_x^2+P_z^2)$,
    $I[5,2] \mathrel{+}= m(P_x^2+P_y^2)$. Off-diagonal products of inertia are not
    added. The default $L_\text{pivot} = 0$ (pivot through origin) gives no correction.
    The grade of the pivot entity -- point in $2D$, line in $3D$ -- matches the grade of
    the velocity generator $\Omega$ (see Tables~\ref{tab:basis_pga2dp}
    and~\ref{tab:basis_pga3dp}).
\end{itemize}

After computing $\dot{\Omega}_b = I_\text{pivot}^{-1}[F_b - \mathrm{rcmt}(\Omega_b,
I_\text{pivot}[\Omega_b])]$ a constraint projection is still required to prevent numerical
drift of off-constraint components during RK4 integration.
In $2D$ the body-frame state is encoded as $B_b = \varphi \cdot Q_b$
and $\Omega_b = \omega \cdot Q_b$; since $Q_b.z = 1$ the scalar angular acceleration is
\begin{equation}
    \alpha = (\dot{\Omega}_b)_z, \qquad \dot{\Omega}_b = \alpha \cdot Q_b
    \label{eq:constraint_projection_2d}
\end{equation}
In $3D$ the analogous projection onto the normalized pivot line gives $\alpha =
\dot{\Omega}_b \cdot \hat{L}_\text{pivot}$. The body-frame encoding $B_b = \varphi \cdot
Q_b$ carries a key geometric property: the world image of $Q_b$ is stationary for all
$\varphi$.
From the body-frame motor
formula~(\ref{eq:motor_pga_body_frame}), $M(t) = M_0 \veedot R(\varphi)$ with the pure
pivot rotation $R(\varphi) = \exp_\veedot(\tfrac{1}{2} B_b(t)) = \exp_\veedot(\tfrac{1}{2}\,
\varphi\, Q_b)$ and the constant base motor $M_0 = M(t{=}0)$. A finite point is invariant
under a rotation about itself, so $R(\varphi) \veedot Q_b \veedot \rrev{R(\varphi)} = Q_b$,
and the world image of the pivot reduces to
\begin{equation}
    M(t) \veedot Q_b \veedot \rrev{M(t)}
    \;=\; M_0 \veedot Q_b \veedot \rrev{M_0}
    \;=\; Q_{b,\text{world}} = \text{const}
    \label{eq:pivot_invariance}
\end{equation}
No explicit constraint enforcement is required for the \emph{position} of the pivot -- it
is preserved exactly for all $\varphi$ by the motor algebra. Only the velocity direction
(eqn.~(\ref{eq:constraint_projection_2d})) must be enforced at each integration step. \\


\newpage

\subsubsection{Structure of inertia and inertia maps}
\label{subsubsection:structure_inertia_maps}

The inertia map $I[\Omega]$ derived in the preceding subsections is a linear map that
takes the velocity twist $\Omega$ of a rigid body and returns its momentum wrench $P =
I[\Omega]$. Its matrix representation -- shown in
eqns.~(\ref{eq:inertia_map_components_2D}) and (\ref{eq:inertia_map_components_3D}) --
encodes mass, scalar moment of inertia, and the full classical inertia tensor in a single
algebraic object. However, the classical quantities do not sit where standard intuition
might suggest: mass does not appear on the diagonal, and moments of inertia do not appear
in the block connecting forces to accelerations of the same type. \\

The placement of each entry is determined by two ingredients: the choice of basis for the
input and output of the inertia map, and the complement structure of the respective
algebra. The right complement $\rcmpl{(\cdot)}$ maps each basis element to its complement
in the algebra. In $PGA$ the translational and rotational generators live in dual
subspaces -- ideal lines are dual to real lines -- so the mass contribution, which couples
velocity to momentum across dual subspaces, always lands in an off-diagonal block. The
following paragraphs explain this structure first for $2D$, then for $3D$, and concludes
with a unified comparison.

\paragraph{$PGA2DP$ --- structure of the 2D inertia map:}
In $2D$ $PGA$ the velocity twist $\Omega$ is a vector and the momentum wrench $P$ is a
bivector, so the inertia map crosses two different grades:
\begin{equation}
    I_{2D}[\Omega] : \mathrm{Vec2dp} \to \mathrm{BiVec2dp}
\end{equation}
The basis elements and their physical roles are listed in Table~\ref{tab:basis_pga2dp}.

\begin{table}[!ht]
\centering\small
\begin{tabular}{c|p{5.2cm}|p{5.2cm}}
\hline
Index & Vec2dp input (velocity twist $\Omega$) & BiVec2dp output (momentum wrench $P$) \\
\hline
0 & $\bv{1}$ -- $x$-velocity & $\bv{31}$ -- $x$-momentum \\
1 & $\bv{2}$ -- $y$-velocity & $\bv{32}$ -- $y$-momentum \\
2 & $\bv{3}$ -- angular velocity  & $\bv{12}$ -- angular momentum \\
\hline
\end{tabular}
\caption{Basis correspondence for the $2D$ $PGA$ inertia map.}
\label{tab:basis_pga2dp}
\end{table}

For a unitized point $(X_x, X_y, X_z)$ with projective coordinate $X_z = 1$ and mass $m$,
the inertia matrix from eqn.~(\ref{eq:inertia_map_components_2D}) simplifies to
\begin{equation}
    I_{2D} = m
    \begin{bmatrix}
        0    &  1   & -X_y           \\
        -1   &  0   &  X_x           \\
        -X_x & -X_y & X_x^2 + X_y^2
    \end{bmatrix}
    \label{eq:inertia_2D_unitized}
\end{equation}
where rows correspond to output components ($\bv{31}$, $\bv{32}$, $\bv{12}$) and columns
to input components ($\Omega_x$, $\Omega_y$, $\Omega_z$).

\paragraph{Why mass is off-diagonal:} The mass contribution to the inertia map takes the
form
$I_\mathrm{mass}(\Omega) = -m \cdot \rcmpl{\Omega}$, where $\rcmpl{(\cdot)}$ denotes the
complement in $PGA2DP$. The complement of the translational basis vectors crosses indices:
\begin{equation}
    \rcmpl{\bv{1}} = \bv{32}, \qquad \rcmpl{\bv{2}} = -\bv{31}
    \label{eq:dual_pga2dp}
\end{equation}
Applying $I_\mathrm{mass} = -m \cdot \rcmpl{(\cdot)}$ gives the following contributions:
\begin{itemize}
    \item $\Omega_x$ (component along $\bv{1}$, column~0) maps to $-m \cdot \bv{32}$,
    landing at output row~1 $\rightarrow$ $I[1,0] = -m$.
    \item $\Omega_y$ (component along $\bv{2}$, column~1) maps to $-m \cdot (-\bv{31})
    = +m \cdot \bv{31}$, landing at output row~0 $\rightarrow$ $I[0,1] = +m$.
\end{itemize}
The diagonal entries $I[0,0]$ and $I[1,1]$ are zero because the complement of $\bv{1}$ is
$\bv{32}$ (not $\bv{31}$) and the complement of $\bv{2}$ is $-\bv{31}$ (not $\bv{32}$):
there is no self-coupling. This is a direct consequence of the degenerate metric $\bv{3}^2
= 0$. \\

The resulting mass block is antisymmetric with $I[0,1] = +m$ and $I[1,0] = -m$. The sign
structure is consistent with the kinematic encoding: the motor-generator $\Omega$ for a
pure translation with classical velocity $(v_x, v_y)$ is encoded in $2D$ $PGA$ as $\Omega
= (-v_y, v_x, 0)$. The two sign flips cancel and the classical linear momentum $p_x =
mv_x$, $p_y = mv_y$ is correctly recovered. In any other context -- position, force,
acceleration -- $\bv{1}$ always encodes the $x$-component and $\bv{2}$ the $y$-component.

\paragraph{Moment of inertia and cross-coupling terms:} The entry $I[2,2] = m(X_x^2 +
X_y^2) =
mr^2$ maps angular velocity ($\bv{3}$, index~2) to angular momentum ($\bv{12}$, index~2),
recovering the classical scalar moment of inertia $J = mr^2$ for a point mass at distance
$r$ from the rotation center. The complement of $\bv{3}$ is $\rcmpl{\bv{3}} = -\bv{12}$,
which connects the angular velocity generator directly to the angular momentum carrier;
the sign is absorbed into the definition so that $J > 0$. The entries $I[0,2] = -mX_y$ and
$I[1,2] = +mX_x$ couple angular velocity to linear momentum, encoding the classical
relation $\vec{L}_\mathrm{linear} = m(\vec{\omega} \times \vec{r})$ for a body rotating at
offset $(X_x, X_y)$ from the rotation center. These entries, together with their
transposes $I[2,0]$ and $I[2,1]$, vanish when the point is at the origin.

\paragraph{Body-frame inertia matrix:} Integrating the single-point formula over a rigid
body
with center of mass at the origin (body frame) gives:
\begin{equation}
    I_\mathrm{body} =
    \begin{bmatrix}
        0  &  m  &  0 \\
        -m &  0  &  0 \\
        0  &  0  &  J
    \end{bmatrix}
    \label{eq:inertia_2D_body}
\end{equation}
where $J = \int_m (x^2 + y^2)\,dm$ is the classical scalar moment of inertia about the
center of mass. The antisymmetric mass block (upper-left $2\times2$) and the moment of
inertia entry (lower-right $1\times1$) are cleanly separated: translational and rotational
degrees of freedom decouple at the center of mass, exactly as in classical mechanics.

\paragraph{General frame with center of mass at $(c_x, c_y)$:} When the chosen reference
point
does not coincide with the center of mass, the integrated inertia matrix is:
\begin{equation}
    I =
    \begin{bmatrix}
        0      &  m     & -m c_y                              \\
        -m     &  0     &  m c_x                              \\
        -m c_x & -m c_y & J_\mathrm{cm} + m(c_x^2 + c_y^2)
    \end{bmatrix}
    \label{eq:inertia_2D_general}
\end{equation}
where $J_\mathrm{cm}$ is the scalar moment of inertia about the center of mass and $r^2 =
c_x^2 + c_y^2$ is the squared distance from the center of mass to the reference point. The
matrix decomposes into the following sub-blocks:
\begin{itemize}
    \item \emph{Upper-left $2\times2$} $[0, m;\; -m, 0]$: the antisymmetric mass block.
    It is \emph{frame-invariant} -- its values are unchanged regardless of where the
    center of mass is located.
    \item \emph{Upper-right $2\times1$} $[-mc_y;\; mc_x]$: angular velocity $\to$ linear
    momentum coupling. Encodes $\vec{L}_\mathrm{linear} = m(\vec{\omega} \times
    \vec{r}_\mathrm{cm})$. Zero only when the center of mass coincides with the reference
    point.
    \item \emph{Lower-left $1\times2$} $[-mc_x,\; -mc_y]$: translational velocity $\to$
    angular momentum coupling. Transpose of the upper-right block (up to sign convention).
    \item \emph{Lower-right $1\times1$}: $J_\mathrm{cm} + m(c_x^2 + c_y^2)$ -- the
    classical parallel axis theorem applied to the scalar moment of inertia.
\end{itemize}

\paragraph{Uniform rectangular plate:} As a verification, integrating over a uniform plate
of
width $W$ (along $\bv{1}$), height $H$ (along $\bv{2}$), total mass $M$, centered at the
origin gives:
\begin{equation}
    I_\mathrm{plate} = M
    \begin{bmatrix}
        0  &  1  &  0 \\
        -1 &  0  &  0 \\
        0  &  0  & \dfrac{W^2 + H^2}{12}
    \end{bmatrix}
    \label{eq:inertia_plate}
\end{equation}
The entry $I[2,2] = M(W^2+H^2)/12$ is the standard formula for the moment of inertia of a
uniform rectangle about its centroid axis. For a discrete endpoint-inclusive grid of $n
\times n$ equally spaced point masses, the relative error in $I[2,2]$ compared to the
continuous value is $2/(n-1)$, independent of plate dimensions, converging from above as
$n \to \infty$. \\

Table~\ref{tab:classical_pga2dp} summarizes the classical inertial quantities and their
locations in the $2D$ $PGA$ inertia matrix.

\begin{table}[h!]
\centering\small
\begin{tabular}{p{4cm}|p{5.5cm}|p{2.8cm}}
\hline
Classical quantity & $2D$ $PGA$ entry & Condition \\
\hline
Total mass $m$ & $I[0,1] = +m$,\; $I[1,0] = -m$ & Always \\
Scalar moment $J$ about CoM & $I[2,2] = J$ & CoM at origin \\
CoM offset $(c_x, c_y)$ & $I[0,2] = -mc_y$,\; $I[1,2] = +mc_x$ (and transposed) & General
frame \\
Parallel axis correction & $I[2,2] = J_\mathrm{cm} + m(c_x^2 + c_y^2)$ & General frame \\
\hline
\end{tabular}
\caption{Classical inertial quantities and their location in the $2D$ $PGA$ $3\times3$
inertia matrix.}
\label{tab:classical_pga2dp}
\end{table}

\paragraph{$3D$ $PGA$ --- structure of the 3D inertia map:}
In $3D$ $PGA$ both the velocity twist $\Omega$ and the momentum wrench $P$ are bivectors,
so the inertia map acts within the same grade:
\begin{equation}
    I_{3D}[\Omega] : \mathrm{BiVec3dp} \to \mathrm{BiVec3dp}
\end{equation}
The six basis bivectors split into two groups: \emph{ideal lines} ($\bv{41}$, $\bv{42}$,
$\bv{43}$, containing the projective direction $\bv{4}$ which squares to zero) carrying
translational information, and \emph{real lines} ($\bv{23}$, $\bv{31}$, $\bv{12}$) of the
Euclidean subspace carrying rotational information. Their physical roles are listed in
Table~\ref{tab:basis_pga3dp}.

\begin{table}[h!]
\centering\small
\begin{tabular}{c|p{5.7cm}|p{5.7cm}}
\hline
Idx & BiVec3dp input (velocity twist $\Omega$) & BiVec3dp output (momentum wrench $P$) \\
\hline
0 & $\bv{41}$ -- ideal line, $x$-velocity &
$\bv{41}$ -- linear momentum $x$
\\
1 & $\bv{42}$ -- ideal line, $y$-velocity &
$\bv{42}$ -- linear momentum $y$
\\
2 & $\bv{43}$ -- ideal line, $z$-velocity &
$\bv{43}$ -- linear momentum $z$
\\
3 & $\bv{23}$ -- real line, angular velocity about $\bv{1}$ &
$\bv{23}$ -- angular momentum about $\bv{1}$
\\
4 & $\bv{31}$ -- real line, angular velocity about $\bv{2}$ &
$\bv{31}$ -- angular momentum about $\bv{2}$
\\
5 & $\bv{12}$ -- real line, angular velocity about $\bv{3}$ &
$\bv{12}$ -- angular momentum about $\bv{3}$
\\
\hline
\end{tabular}
\caption{Basis correspondence for the $3D$ $PGA$ inertia map.}
\label{tab:basis_pga3dp}
\end{table}

The $6\times6$ inertia matrix from eqn.~(\ref{eq:inertia_map_components_3D}) for a
unitized point with $X_w = 1$ is:
\begin{align}
    I_{3D} = m
    \begin{bmatrix}
      0              &  X_z           & -X_y           & 1    &  0   &  0   \\
      -X_z           &  0             &  X_x           & 0    &  1   &  0   \\
       X_y           & -X_x           &  0             & 0    &  0   &  1   \\
       X_y^2 + X_z^2 & -X_x X_y       & -X_x X_z       & 0    & -X_z &  X_y \\
      -X_x X_y       &  X_x^2 + X_z^2 & -X_y X_z       & X_z  &  0   & -X_x \\
      -X_x X_z       & -X_y X_z       &  X_x^2 + X_y^2 & -X_y &  X_x &  0
    \end{bmatrix}
    \label{eq:inertia_3D_unitized}
\end{align}
where columns correspond to ($\Omega_{vx}$, $\Omega_{vy}$, $\Omega_{vz}$, $\Omega_{mx}$,
$\Omega_{my}$, $\Omega_{mz}$) and rows to the same ordering for $P$.

\paragraph{Block structure and the complement in 3D:} The $6\times6$ matrix decomposes
into
four $3\times3$ blocks with distinct physical roles:
\begin{itemize}
    \item \emph{Upper-left [0:3, 0:3]}: translational velocity $\to$ linear momentum. For
    a point at the origin this block is zero. For a point displaced to $(X_x, X_y, X_z)$
    it becomes the antisymmetric cross-product matrix $[X\times]$, encoding the linear
    momentum contribution due to the off-axis position.
    \item \emph{Upper-right [0:3, 3:6]}: angular velocity $\to$ linear momentum. For a
    unitized point ($X_w = 1$) this block equals $m\,\mathbf{I}_3$ (mass times the
    $3\times3$ identity). This is the $3D$ analogue of the off-diagonal mass entries in
    the $2D$ case.
    \item \emph{Lower-left [3:6, 0:3]}: translational velocity $\to$ angular momentum.
    This block carries the classical moments and products of inertia.
    \item \emph{Lower-right [3:6, 3:6]}: angular velocity $\to$ angular momentum. Zero for
    a point at the origin; becomes antisymmetric for a displaced center of mass.
\end{itemize}

The placement of mass in the upper-right block follows from the complement in $3D$ $PGA$,
which maps real lines to ideal lines:
\begin{equation}
    \rcmpl{\bv{23}} = \bv{41}, \quad \rcmpl{\bv{31}} = \bv{42}, \quad
    \rcmpl{\bv{12}} = \bv{43}
    \label{eq:dual_pga3dp}
\end{equation}
Applying $I_\mathrm{mass} = -m \cdot \rcmpl{(\cdot)}$ maps each angular velocity component
to the corresponding linear momentum component, placing $m$ on the diagonal of the
upper-right block: $I[0,3] = I[1,4] = I[2,5] = m$. The diagonal entries of the upper-left
and lower-right blocks are zero by the same degenerate-metric argument as in $2D$.

\paragraph{Body-frame inertia matrix:} Integrating the single-point formula over a rigid
body
with center of mass at the origin gives:
\begin{equation}
    I_\mathrm{body} =
    \begin{bmatrix}
        0       & 0       & 0       &  m   & 0    & 0    \\
        0       & 0       & 0       &  0   & m    & 0    \\
        0       & 0       & 0       &  0   & 0    & m    \\
        I_{xx}  & -I_{xy} & -I_{xz} &  0   & 0    & 0    \\
        -I_{xy} &  I_{yy} & -I_{yz} &  0   & 0    & 0    \\
        -I_{xz} & -I_{yz} &  I_{zz} &  0   & 0    & 0
    \end{bmatrix}
    \label{eq:inertia_3D_body}
\end{equation}
written in block form as $I_\mathrm{body} = [\mathbf{0},\; m\mathbf{I}_3;\;
\mathbf{J}_{cm},\; \mathbf{0}]$, where $\mathbf{J}_{cm}$ is the classical $3\times3$
inertia tensor about the center of mass:
\begin{equation}
    \mathbf{J}_{cm} =
    \begin{bmatrix}
        I_{xx}  & -I_{xy} & -I_{xz} \\
        -I_{xy} &  I_{yy} & -I_{yz} \\
        -I_{xz} & -I_{yz} &  I_{zz}
    \end{bmatrix}
    \label{eq:classical_inertia_tensor}
\end{equation}
with the classical integrals for the diagonal moments $I_{xx} = \int_m (y^2+z^2)\,dm$,
$I_{yy} = \int_m (x^2+z^2)\,dm$, $I_{zz} = \int_m (x^2+y^2)\,dm$ and the off-diagonal
products of inertia $I_{xy} = \int_m xy\,dm$, $I_{xz} = \int_m xz\,dm$, $I_{yz} = \int_m
yz\,dm$. Translational and rotational degrees of freedom decouple at the center of mass:
the diagonal blocks are zero and all physical content sits in the off-diagonal blocks.

\paragraph{General frame with center of mass at $\vec{r}_{cm} = (c_x, c_y, c_z)$:} In a
general
frame all four $3\times3$ blocks are non-zero:
\begin{equation}
    I =
    \begin{bmatrix}
        0       &  m c_z   & -m c_y   &  m   &  0     &  0     \\
        -m c_z  &  0       &  m c_x   &  0   &  m     &  0     \\
         m c_y  & -m c_x   &  0       &  0   &  0     &  m     \\
        I_{xx}  & -I_{xy}  & -I_{xz}  &  0   & -m c_z &  m c_y \\
        -I_{xy} &  I_{yy}  & -I_{yz}  &  m c_z &  0   & -m c_x \\
        -I_{xz} & -I_{yz}  &  I_{zz}  & -m c_y &  m c_x &  0
    \end{bmatrix}
    \label{eq:inertia_3D_general}
\end{equation}
where the lower-left inertia tensor entries include the classical parallel axis
corrections:
\begin{align}
    I_{xx} &= I_{xx,cm} + m(c_y^2 + c_z^2), \quad &I_{xy} &
    = I_{xy,cm} + m c_x c_y \nonumber \\
    I_{yy} &= I_{yy,cm} + m(c_x^2 + c_z^2), \quad &I_{xz} &
    = I_{xz,cm} + m c_x c_z \label{eq:parallel_axis_3D} \\
    I_{zz} &= I_{zz,cm} + m(c_x^2 + c_y^2), \quad &I_{yz} &
    = I_{yz,cm} + m c_y c_z \nonumber
\end{align}
Two structural features carry over from the $2D$ case:
\begin{itemize}
    \item \emph{The upper-right block $m\,\mathbf{I}_3$ is frame-invariant.} It remains
    equal to $m\,\mathbf{I}_3$ regardless of the center-of-mass position. Mass is a dual
    coupling between ideal and real lines that does not depend on position. This is the
    $3D$ counterpart of the invariant antisymmetric mass block in $2D$.
    \item \emph{The upper-left and lower-right blocks are conjugate antisymmetric
    matrices.} Upper-left~$= m[\vec{r}_{cm}\times]^\top$ and lower-right~$=
    m[\vec{r}_{cm}\times]$, where $[\vec{r}_{cm}\times]$ is the $3\times3$ antisymmetric
    cross-product matrix of the center-of-mass position vector. Both blocks are driven by
    the same offset $\vec{r}_{cm}$, reflecting that the translational-to-translational and
    angular-to-angular couplings induced by the offset are conjugate couplings.
\end{itemize}

This $6\times6$ structure is the $PGA$ form of the classical spatial inertia matrix from
rigid body dynamics~\cite{Dynamics_in_PGA_plane-based_Dorst-DeKennik:2023}, expressed in
the (translational, angular) velocity ordering of $3D$ $PGA$.

\paragraph{Uniform rectangular cuboid:} For a uniform cuboid of width $w$ (along $\bv{1}$),
height $h$ (along $\bv{2}$), depth $d$ (along $\bv{3}$), total mass $m$, centered at the
origin, all cross-inertia and center-of-mass coupling terms vanish:
\begin{equation}
    I_\mathrm{cuboid} = m
    \begin{bmatrix}
        0 & 0 & 0 & 1 & 0 & 0 \\[4pt]
        0 & 0 & 0 & 0 & 1 & 0 \\[4pt]
        0 & 0 & 0 & 0 & 0 & 1 \\[4pt]
        \dfrac{h^2+d^2}{12} & 0 & 0 & 0 & 0 & 0 \\[10pt]
        0 & \dfrac{w^2+d^2}{12} & 0 & 0 & 0 & 0 \\[10pt]
        0 & 0 & \dfrac{w^2+h^2}{12} & 0 & 0 & 0
    \end{bmatrix}
    \label{eq:inertia_cuboid}
\end{equation}
The three diagonal entries of the lower-left block are the classical rectangular moments
of inertia: $I[3,0] = m(h^2+d^2)/12$ about the $x$-axis, $I[4,1] = m(w^2+d^2)/12$ about
the $y$-axis, and $I[5,2] = m(w^2+h^2)/12$ about the $z$-axis -- each determined by the
two extents of the body perpendicular to the respective rotation axis, as in the classical
formula. \\

Table~\ref{tab:classical_pga3dp} gives the correspondence between classical inertial
quantities and their location in the $3D$ $PGA$ inertia matrix.

\begin{table}[h!]
\centering\small
\begin{tabular}{p{3.1cm}|p{4.6cm}|p{4.7cm}}
\hline
Classical quantity & $3D$ $PGA$ block & Entries (body frame, CoM at origin) \\
\hline
Total mass $m$ & Upper-right diagonal ($m\mathbf{I}_3$) & $I[0,3]=I[1,4]=I[2,5]=m$ \\
Diagonal moments $I_{xx}, I_{yy}, I_{zz}$ & Lower-left diagonal & $I[3,0],\; I[4,1],\;
I[5,2]$ \\
Products of inertia $I_{xy}, I_{xz}, I_{yz}$ & Lower-left off-diagonal &
$I[3,1]=I[4,0]=-I_{xy}$, etc. \\
CoM offset (general frame) & All four $3\times3$ blocks & via parallel axis theorem \\
\hline
\end{tabular}
\caption{Classical inertial quantities and their location in the $3D$ $PGA$ $6\times6$
inertia matrix (body frame, center of mass at origin).}
\label{tab:classical_pga3dp}
\end{table}

\paragraph{Unified view --- 2D and 3D:}
In both $2D$ $PGA$ and $3D$ $PGA$ the mass contribution to the inertia map takes the same
algebraic form:
\begin{equation}
    I_\mathrm{mass}(\Omega) = -m \cdot \rcmpl{\Omega}
    \label{eq:inertia_mass_dual}
\end{equation}
The complement maps each velocity generator to its momentum carrier in the dual subspace,
placing mass in the off-diagonal block that connects dual velocity and momentum spaces.
The location of that block in the matrix depends on the basis ordering. \\

In $2D$ $PGA$ the translational generators ($\bv{1}$, $\bv{2}$, indices 0--1) sit at lower
indices and the rotational generator ($\bv{3}$, index~2) sits at the highest index. The
complement maps translational inputs to translational outputs while crossing $x$ and $y$,
placing mass in the \emph{upper-left $2\times2$} sub-block (translational velocity $\to$
linear momentum). \\

In $3D$ $PGA$ the translational generators ($\bv{41}$, $\bv{42}$, $\bv{43}$, indices 0--2)
are ideal lines at lower indices, while the rotational generators ($\bv{23}$, $\bv{31}$,
$\bv{12}$, indices 3--5) are real lines at higher indices. The complement maps real-line
inputs (angular velocity) to ideal-line outputs (linear momentum), placing mass in the
\emph{upper-right $3\times3$} sub-block (angular velocity $\to$ linear momentum). \\

Table~\ref{tab:unified_pga} gives the side-by-side comparison for the body frame (center
of mass at origin).

\begin{table}[h!]
\centering\small
\begin{tabular}{p{3.1cm}|p{4.1cm}|p{5.1cm}}
\hline
Classical quantity & $2D$ $PGA$ ($3\times3$) & $3D$ $PGA$ ($6\times6$) \\
\hline
Total mass $m$ &
    Upper-left $2\times2$: $I[0,1]=+m$, $I[1,0]=-m$ &
    Upper-right $3\times3$ diag.: $I[0,3]=I[1,4]=I[2,5]=m$ \\
Scalar/diagonal moments &
    Lower-right $1\times1$: $I[2,2]=J$ &
    Lower-left $3\times3$ diag.: $I[3,0], I[4,1], I[5,2]$ \\
Products of inertia &
    n/a (scalar $J$ in $2D$) &
    Lower-left $3\times3$ off-diagonal \\
CoM coupling &
    $I[0,2], I[1,2]$ and transposed &
    All four $3\times3$ blocks non-zero \\
Parallel axis &
    $I[2,2] = J_\mathrm{cm} + mr^2$ &
    Lower-left $3\times3$ (3D formula) \\
\hline
\end{tabular}
\caption{Unified mapping of classical inertial quantities to the $PGA$ inertia matrix
(body frame, center of mass at origin).}
\label{tab:unified_pga}
\end{table}

The frame-invariance of the mass block and the decoupling at the center of mass are common
to both dimensions. In the body frame with the origin at the center of mass -- the natural
choice for integrating eqn.~(\ref{eq:ode_system_body_frame}) -- the coupling blocks
vanish, the mass block takes its simplest form, and the moment block reduces to $J$ (in
$2D$) or $\mathbf{J}_{cm}$ (in $3D$) without any parallel axis corrections. This matches
the classical result that Newton's and Euler's equations decouple at the center of mass,
and it is the reason the body frame with the center of mass at the origin is used as the
preferred reference frame throughout this section. Note, however, that while the inertia
map $I_\mathrm{body}$ is block-diagonal at the center of mass -- translational and
rotational blocks decoupling in the inertia matrix itself -- the equations of motion
retain dynamical coupling through the $\mathrm{rcmt}(\Omega_b, I_b[\Omega_b])$ term. This
is the $PGA$ counterpart of the classical coupling terms $\vec{\omega}\times\vec{v}$ and
$\vec{\omega}\times\mathbf{I}_b\vec{\omega}$, and it arises from the rotation of the body
frame, not from the inertia structure.

\paragraph{Pivot-corrected inertia --- using the pivot as reference point:}
The general-frame formulae derived above -- eqn.~(\ref{eq:inertia_2D_general}) in $2D$ and
eqn.~(\ref{eq:inertia_3D_general}) in $3D$ -- apply directly when the reference point is
chosen to coincide with the pivot rather than the center of mass. In that case the
center-of-mass offset vector $(c_x, c_y)$ resp.\ $\vec{r}_{cm}$ points from the pivot to
the center of mass, and the inertia map automatically incorporates the full Steiner
correction: $I[2,2] = J_{cm} + mr^2$ in $2D$ (eqn.~(\ref{eq:steiner_2d})) and the complete
tensor update of eqn.~(\ref{eq:parallel_axis_3D}) in $3D$.

\newpage

\subsubsection{Moving coordinate systems}
\label{subsubsection:moving_coordinate_systems}

\begin{figure}[H]
    \includegraphics[width=7cm]{moving_frame.png}
    \centering
    \caption{Moving frame: World or reference frame \{$X,Y$\} (only $2D$ shown as
    examples, extends naturally to $3D$). Body frame \{$X_b, Y_b$\}. Point $A$ might move
    with $v_A$ and accelerate with $a_A$ (not shown). Body frame turns with rotational
    speed $\omega$ and rotational acceleration $\alpha$ relative to the world frame.}
    \label{fig:moving_frame}
\end{figure}

When modelling motion in moving coordinate systems we look into relative motion.
Typically, we use a world frame as reference frame and a body frame as frame moving
relative to the world frame, as shown in fig.~\ref{fig:moving_frame}. $x_A$ and $x_B$ are
absolute positions in the world frame. \\

Most of the time $x_A$ is a reference position (e.g. the center of mass of a rigid
body)\footnote{The reference position $x_A$ can be any arbitrary point of the rigid body.
However, it is often convenient to choose $x_A$ to be located at the center of mass.}, and
we are interested in the movement of $x_B$ relative to $x_A$. The relative position of $B$
with respect to $A$ is shown as $r_{B/A}$. If A is the center of mass of the rigid body
then $r$ are local (relative) coordinates of the points of the body in the body frame. So
we define
\begin{subeqnarray}
    r        &=& x_B - x_A \quad \textit{(relative position)} \slabel{eq:rel_pos}\\
    \dot{r}  &=& v_B - v_A \quad \textit{(relative velocity)} \slabel{eq:rel_spd}\\  
    \ddot{r} &=& a_B - a_A \quad \textit{(relative acceleration)} \slabel{eq:rel_acl}
\end{subeqnarray}
as relative position, speed and acceleration. \\

The relative position in a $3D$ frame is defined with $(x,y,z)$ as local or relative
coordinates
\begin{equation}
    r = x \bv{1} + y \bv{2} + z \bv{3}
\end{equation}
The total change of the relative position has the following components
\begin{equation}
    \dot{r} = \dfdt{r}{t} = \underbrace{\dot{x} \bv{1} + \dot{y} \bv{2} + \dot{z} \bv{3}}
              _{\textrm{changing speed}: \; \dfdt{r}{t} \;=\; v}
            + \underbrace{x \dfdt{\bv{1}}{t} + y \dfdt{\bv{2}}{t} + z \dfdt{\bv{3}}{t}}
              _{\textrm{changing frame}: \; \dfdt{r}{t} \;=\; \omega \times r}
\end{equation}
leading to the relative speed $\dot{r}$
\begin{equation}
    \dot{r} = \omega \times r_{B/A} + v_{B/A}
    \label{eq:rel_spd_two_terms}
\end{equation}
which can be transformed using eqn.~(\ref{eq:rel_spd}) resulting to the total speed at
position $B$ as
\begin{equation}
    v_B = v_A + \omega \times r_{B/A} + v_{B/A} 
\end{equation}
We can get the relative acceleration from deriving eqn.~(\ref{eq:rel_spd_two_terms}) again
with respect to $t$
\begin{equation}
    \ddot{r} = \dot{\omega} \times r_{B/A} + \omega \times \dot{r}_{B/A} + \dot{v}_{B/A}
    \label{eq:rel_acl_three_terms}
\end{equation}
Using eqn.~(\ref{eq:rel_spd_two_terms}) directly for $\dot{r}_{B/A}$ and for
$\dot{v}_{B/A} = \omega \times v_{B/A} + a_{B/A}$ this yields
\begin{equation}
    \ddot{r} = \dot{\omega} \times r_{B/A} + \omega \times \omega \times r_{B/A}
             + \omega \times v_{B/A} + \omega \times v_{B/A} + a_{B/A}
    \label{eq:rel_acl_five_terms}
\end{equation}
when reordered for terms depending on position, relative speed and relative acceleration,
we get
\begin{equation}
    \ddot{r} = \underbrace{\omega \times \omega \times r_{B/A}}_
                          {\textrm{centripetal accel.}
                          \;=\; -|\omega|^2 \; r_{B/A}}
             + \underbrace{2\; \omega \times v_{B/A}}_
                          {\textrm{Coriolis acceleration}}
             + \underbrace{\dot{\omega} \times r_{B/A}}_
                          {\textrm{frame acceleration}}
             + \underbrace{a_{B/A}}_
                          {\textrm{particle accel.}}
    \label{eq:rel_acl_four_terms}
\end{equation}
Using eqn.~(\ref{eq:rel_acl}) this results in the total acceleration at position $B$ as
\begin{equation}
    a_B = a_A + \omega \times \omega \times r_{B/A}
              + 2\; \omega \times v_{B/A}
              + \dot{\omega} \times r_{B/A} + a_{B/A}
\end{equation}

Now we can translate the various velocity and acceleration terms that describe physics
using expressions in the geometric algebra language, separately for $EGA$ and $PGA$. To
get there, we mainly replace the cross-product expressions by bivectors in $EGA$ and
trivectors in $PGA$. They are duals of the momentary rotation axis. This directly applies
to $2D$ in that form.\footnote{The $2D$ application is done in classical vector algebra by
using the $z$-component of the $3D$ cross-product as replacement, even if the
cross-product is not defined for $2D$. In the classical approach we temporarily have to
include the third dimension and use the relevant $z$-component exclusively. However, this
temporary inclusion of an additional dimension is not required in geometric algebra, since
a rotation only requires two dimensions for its representation in geometric algebra.}

% Three-column layout: Application (spanning) | Tag | Description
% Adjust \cwAppPhys and \cwTagPhys; description fills the rest.
\newlength{\cwAppPhys}\setlength{\cwAppPhys}{0.15\textwidth}
\newlength{\cwTagPhys}\setlength{\cwTagPhys}{0.06\textwidth}
\newlength{\cwDescPhys}\setlength{\cwDescPhys}{\dimexpr\textwidth - \cwAppPhys -
\cwTagPhys - 4\arrayrulewidth - 6\tabcolsep\relax}
\begin{table}[H]
{\setlength{\extrarowheight}{2pt}\renewcommand{\arraystretch}{1.15}%
% \extrarowheight: absolute extra space added above each row (before top rule);
% \arraystretch:   scales the row strut height (1.15 ≈ +15%, adds ~1.5pt below content).
% Together: ~2pt top padding, ~1.5pt bottom padding per row.
\begin{tabular}{|p{\cwAppPhys}|p{\cwTagPhys}|p{\cwDescPhys}|}
    \hline
    \textbf{Term} & & \textbf{Description} \\
    \hline \hline
    \multirow{3}{\cwAppPhys}{Definition of vectors $v$ and points $P$}
    & VA  & $v=(v_x, v_y, v_z)$, $P=(p_x, p_y, p_z)$ \\
    \cline{2-3}
    & EGA & $v=(v_x, v_y, v_z)$, $P=(p_x, p_y, p_z)$ \\
    \cline{2-3}
    & PGA & $v=(v_x, v_y, v_z, 0)$, $P=(p_x, p_y, p_z, 1)$ \\
    \hline \hline
    \multirow{3}{\cwAppPhys}{Definition of rotation}
    & VA  & rotation axis vector $\omega = (\omega_x,\omega_y,\omega_z)$\newline
    (a direction, axis always contains origin)\\
    \cline{2-3}
    & EGA & rotation bivector $\Omega_{E} = \rightbulkdual{\omega} =
    (\Omega_x,\Omega_y,\Omega_z)$\newline
    (bivector dual to direction vector $\omega$ -- defines the plane of rotation, which
    always contains the origin)\\
    \cline{2-3}
    & PGA & rotation direction $\omega = (\omega_x,\omega_y,\omega_z, 0)$ \newline
    reference point on rotation axis $Q = (q_x,q_y,q_z, 1)$
    rotation bivector $\Omega_{P} = Q \wedge \omega $ - a line, where
    \newline
    $\Omega_{P} = (\Omega_{vx},\Omega_{vy},\Omega_{vz},
    \Omega_{mx},\Omega_{my},\Omega_{mz})$
    \newline
    (a rotation axis modelled by a bivector, with arbitrary orientation \emph{and}
    location; can also be created from any unitized line that is scaled with |$\omega$|)\\
    \hline \hline
    \multirow{3}{\cwAppPhys}{Speed of rotation at $r$ or $R$\newline}
    & VA  & evaluation of velocity at $r = (r_x,r_y,r_z)$
    ($r$ is vector perpendicular to axis; points to evaluation point)\newline
    $v = \omega \times r$  \\
    \cline{2-3}
    & EGA & evaluation of velocity at $r = (r_x,r_y,r_z)$
    (see VA)\newline
    $v = \textrm{cmt}(r, \Omega_{E}) = \Omega_{E} \gg r$ \\
    \cline{2-3}
    & PGA & evaluation of velocity at $R = (r_x,r_y,r_z, 1)$\newline
    $v = \textrm{rcmt}(\Omega_{P}, R)$  \\
    \hline \hline
    \multirow{3}{\cwAppPhys}{Centripetal acceleration}
    & VA  & $a_{cp} = \omega \times \omega \times r_{B/A}$ \\
    \cline{2-3}
    & EGA & $a_{cp} = \textrm{cmt}(\textrm{cmt}(r_{B/A}, \Omega_{E}), \Omega_{E})
     = \Omega_{E} \gg (\Omega_{E} \gg r_{B/A})$ \\
    \cline{2-3}
    & PGA & $a_{cp} = \textrm{rcmt}(\Omega_{P}, \textrm{rcmt}(\Omega_{P},r_{B/A}))$
    \\
    \hline \hline
    \multirow{3}{\cwAppPhys}{Coriolis acceleration}
    & VA  & $a_{co} = 2\; \omega \times v_{B/A}$ \\
    \cline{2-3}
    & EGA & $a_{co} = 2\; \textrm{cmt}(v_{B/A}, \Omega_{E})
     = 2 (\Omega_{E} \gg v_{B/A})$ \\
    \cline{2-3}
    & PGA & $a_{co} = 2\; \textrm{rcmt}(\Omega_{P},v_{B/A})$ \\
    \hline \hline
    \multirow{3}{\cwAppPhys}{Frame rotational acceleration}
    & VA  & $a_{frame \; rotation} = \dot{\omega} \times r_{B/A}$ \\
    \cline{2-3}
    & EGA & $a_{frame \; rotation} = \textrm{cmt}(r_{B/A}, \dot{\Omega}_{E})
     = (\dot{\Omega}_{E} \gg r_{B/A})$ \\
    \cline{2-3}
    & PGA & $a_{frame \; rotation} = \textrm{rcmt}(\dot{\Omega}_{P},r_{B/A})$ \\
    % \hline \hline
    % \multirow{3}{\cwAppPhys}{...}
    % & VA  & ... \\
    % \cline{2-3}
    % & EGA & ... \\
    % \cline{2-3}
    % & PGA & ... \\
    \hline
\end{tabular}}% end \extrarowheight + \arraystretch group
\caption{Description of rotation and rotation speed evaluation for \textbf{VA} = Vector
         Algebra (classical approach using vectors); \textbf{EGA} = Euclidean Geometric
         Algebra; \textbf{PGA} = Projective (Euclidean) Geometric Algebra.}
\end{table}

% eqn.~(\ref{})


\subsubsection{Frame trees: static, kinematic and dynamic systems}
\label{subsubsection:frame_trees_systems}

The preceding subsections developed the mechanics of a \emph{single} rigid body: its pose
evolves by a motor $M(t)$, its momentum wrench is the inertia map $P = I[\Omega]$
(subsubsection~\ref{subsubsection:structure_inertia_maps}), and its motion obeys the
body-frame Euler equation~(\ref{eq:ode_system_body_frame}). Modelling assemblies -- a
vehicle with wheels, a robot arm, a pendulum chain -- requires composing many such bodies,
each described \emph{relative} to another. The library provides this as a strict
three-layer stack of classes, all sharing one geometry, a \emph{tree of coordinate frames}.
Each layer inherits the one below and adds exactly one new kind of state:

\begin{table}[H]
\centering\small
\begin{tabular}{|p{0.29\textwidth}|p{0.31\textwidth}|p{0.28\textwidth}|}
    \hline
    \textbf{Layer} & \textbf{Adds per frame} & \textbf{Answers question} \\
    \hline \hline
    \path{static_system2dp} & a pose (relative to parent) & \emph{Where} is each frame? \\
    \hline
    \path{kinematic_system2dp} & a velocity + acceleration twist & \emph{How} is it
    moving? (no forces) \\
    \hline
    \path{dynamic_system2dp} & inertia + joints + gravity & How do \emph{forces} move it?
    \\
    \hline
\end{tabular}
\caption{The three-layer system stack. Each layer inherits the previous one and adds a
         single kind of per-frame state.}
\end{table}

The layering keeps each capability additive and independently testable: the transformation
algebra at the static layer, velocity and acceleration transport at the kinematic layer,
and energy conservation at the dynamic layer. A user who only needs coordinate transforms
never pays for the dynamics. The same three classes exist for $3D$
(\path{static_system3dp}, \path{kinematic_system3dp}, \path{dynamic_system3dp}); their
class structure, method names and algorithms are identical -- only the representations
change, as summarized at the end of this subsection. The tree describes \emph{open}
chains; closed kinematic loops (parallel mechanisms) are added by a separate layer on top of
it, developed in subsection~\ref{subsubsection:closed_loop}.


\subsubsection{Pose, frame and the frame tree}
\label{subsubsection:pose_frame_tree}

A \emph{pose} is the rigid placement of one frame relative to another: an origin and an
orientation (in $2D$ a position vector and a scalar angle $\varphi$; in $3D$ a position
and an axis-angle vector). It is the human-readable form of a \emph{motor} $M$, and the
two interconvert: the motor is built from the pose as in
eqn.~(\ref{eq:motor_pga_body_frame}), and the pose is read back by decoding the origin
from $M \veedot O \veedot \rrev{M}$ and the orientation from the rotated $\bv{1}$-axis (a
constrained motor logarithm, exact on round-trip).

A \path{static_system2dp} holds a \emph{tree} of frames, each recording its parent, with
the crucial invariant that \textbf{a frame's pose is always stored relative to its
parent}, never in world coordinates. The world (root) frame is its own parent. Adding a
frame without naming a parent links it to the previously added one, so a plain sequence
builds a linear chain $\textit{root} \to f_1 \to f_2 \to \dots$, while naming an earlier
frame branches the tree. This provides three modelling capabilities: \emph{branching}
(a body frame with two arms, each arm chain parented off the body, the body off the
possibly-moving world -- any sub-assembly is a subtree); \emph{name addressing} (frames
carry unique names and every operation has a by-name overload); and \emph{reposing}
(moving a frame relative to its parent drags the whole subtree along automatically,
because all descendant poses are relative).\footnote{The frame tree is exercised
end-to-end in the visualization tool \path{ga_view}: the frame-transformation scene
(\path{ga_view/src/active_frame_trafo}) drives a \path{static_system2dp} directly, and the
merry-go-round scene (\path{ga_view/src/active_merry_go_round}) carries the tree one layer
up as a rotating \path{kinematic_system2dp}. Assertion-checked usage is in the test cases
\path{ga_test/src/ga_appl2dp_physics_test.hpp} and its $3D$ analogue
\path{ga_appl3dp_physics_test.hpp}.}

\paragraph{Tree walk:} The transform from any frame to any other is composed along the
tree path. Walking
\emph{up} from the source to the lowest common ancestor (LCA) composes child-to-parent
steps (the regressive reverse of each frame's pose motor); walking \emph{down} from the
LCA to the target composes parent-to-child steps; the result is the regressive product of
the two segments. Special cases fall out of the same procedure: a frame to itself yields
the identity (the unit pseudoscalar, neutral for $\veedot$); a pure ancestor/descendant
relation makes one segment empty (a forward or backward transform); a linear chain is the
degenerate tree where the LCA is simply the shallower frame. Because composition is motor
multiplication with $\veedot$, the entire tree machinery is dimension-neutral and lifts to
$3D$ unchanged.

\paragraph{One sandwich for all objects:} The transform between two frames is a single
motor $M$, and \emph{any} geometric object is
moved by the \emph{same} sandwich $X' = M \veedot X \veedot \rrev{M}$. Classical
mechanics, by contrast, provides several different rules and the modeller must determine
which one applies:
\begin{itemize}
    \item a \emph{point} (a location) transforms \emph{affinely} -- rotate \emph{and}
          translate;
    \item a \emph{free vector} (a direction, a velocity) transforms \emph{linearly} --
          rotate only, no translation;
    \item \emph{basis vectors} transform covariantly while a vector's \emph{coordinates}
          transform contravariantly (the inverse-transpose), the source of much index
          bookkeeping.
\end{itemize}
In $PGA$ there is a single operation, and the object itself carries the information that
selects the right behaviour through its \emph{homogeneous weight} (the projective
component $\bv{3}$ in $2D$, $\bv{4}$ in $3D$): a point has weight $1$ and is rotated
\emph{and} translated; a direction has weight $0$ (a point at infinity) and is rotated
\emph{only}, since translating a point at infinity does nothing. No
contravariant/covariant distinction remains: the motor sandwich is a single
conjugation that transports every grade (points, lines, directions, and --
see below -- twists and wrenches) correctly. The choice of transformation law is thus
handled by the algebra itself.


\subsubsection{Transport of motion along a kinematic chain}
\label{subsubsection:transport_along_chain}

The preceding subsections transport the twist of a \emph{single} body between its body and
world frames (eqn.~(\ref{eq:body_to_world_pga}), eqn.~(\ref{eq:rcmt_pga_world})). The
\path{kinematic_system2dp} generalizes this to a \emph{chain} of frames: each frame
carries a relative velocity twist and a relative acceleration twist with respect to its
parent, and the layer answers velocity and acceleration queries for any point of any
frame.

\paragraph{Adjoint:} A twist $B$ (a vector in $2D$, a bivector in $3D$) is the
instantaneous motion (Lie-algebra) with the motor $M$ defined as $M =
\exp_\veedot(\tfrac12 B)$ (twists \emph{add}, motors \emph{multiply}; see
subsection~\ref{subsection:motor_generators_PGA}). Re-expressing a twist $B$ in another
frame is done by the \emph{adjoint} of the twist, and the geometric-algebra fact that
makes the implementation compact is that the adjoint \emph{is} the very sandwich used to
move a point (or any other object):
\begin{equation}
    \mathrm{Ad}_M(B) = M \veedot B \veedot \rrev{M}
    \label{eq:adjoint_is_sandwich}
\end{equation}
This is eqn.~(\ref{eq:body_to_world_pga}) read as a twist transport rather than an object
transport. So the same sandwich serves double duty -- point action and twist adjoint --
because both are the conjugation $M\veedot(\cdot)\veedot\rrev{M}$. A separate $6\times6$
adjoint matrix, as the classical formulation requires, is unnecessary.

\paragraph{Velocity composition:} This step goes beyond the single-body case. In a chain,
each frame $k$
contributes a relative twist $\xi_k$ expressed in \emph{its own} frame. Transporting each
to the world frame by the adjoint and superposing -- valid because twists live in a vector
space -- gives the world velocity twist of frame $i$:
\begin{equation}
    V_i = \sum_{k \,\in\, \textit{path}(\textit{root} \to i)} \mathrm{Ad}_{M_k}(\xi_k)
        = \sum_k M_k \veedot \xi_k \veedot \rrev{M_k}
    \label{eq:twist_world_sum}
\end{equation}
They add rather than compose precisely because they are infinitesimal: finite motions do
not commute and must be multiplied, but to first order simultaneous infinitesimal motions
sum. Equation~(\ref{eq:twist_world_sum}) is the geometric-algebra form of the robotics
\emph{spatial Jacobian} -- a quantity a single moving body does not exhibit.

\paragraph{Velocity and acceleration fields:} Given a frame's world velocity twist $V$
(and acceleration twist $A$), the velocity of any point $X$ of that frame is the
regressive commutator $\dot{X} = \mathrm{rcmt}(V, X)$, and its acceleration is
$\mathrm{rcmt}(A, X) + \mathrm{rcmt}(V, \mathrm{rcmt}(V, X))$ -- exactly the frame (Euler)
and centripetal terms of subsection~\ref{subsubsection:moving_coordinate_systems}, now
evaluated per chain frame. Differentiating the adjoint along the moving chain yields the
recursive Newton-Euler propagation
\begin{subeqnarray}
    V_i &=& V_{\textit{parent}} + \mathrm{Ad}_{M_i}(\xi_i) \\
    A_i &=& A_{\textit{parent}} + \mathrm{Ad}_{M_i}(\dot{\xi}_i)
            + \big[\, V_i,\; \mathrm{Ad}_{M_i}(\xi_i) \,\big]
    \label{eq:newton_euler_recursion}
\end{subeqnarray}
where the second acceleration term -- the Lie bracket of twists -- is the
Coriolis/centrifugal coupling, measuring how far two twists fail to commute. In this
primal $PGA$ that bracket is again a native product, $[A,B] = \mathrm{rcmt}(A,B)$ (the
$\mathfrak {se}(2)/\mathfrak{se}(3)$ twist bracket). The single-frame transport theorem is
the content of subsection~\ref{subsubsection:moving_coordinate_systems}; the recursion
over a chain is what the kinematic layer adds.

The velocity/finite distinction underlying this -- twists add in the Lie algebra while
motors multiply non-commutatively, equal only when the generators commute (Chasles'
theorem, the screw) -- is developed in
subsection~\ref{subsubsection_lie_algebra_vs_lie_group}. It is also the reason the
parent-to-child pose of a frame is built by \emph{composing} a rotation and a translation
motor rather than exponentiating their summed generators.


\subsubsection{Articulated coupled dynamics}
\label{subsubsection:articulated_dynamics}

The earlier subsections solve a \emph{single} free rigid body
(eqn.~(\ref{eq:ode_system_body_frame})) and a \emph{single} revolute pivot
(eqn.~(\ref{eq:pivot_invariance}), eqn.~(\ref{eq:constraint_projection_2d})). The
\path{dynamic_system2dp} adds the coupled dynamics: \emph{several bodies coupled by
joints}, integrated together in \emph{reduced (joint) coordinates}.\footnote{The canonical
example is the double pendulum -- two revolute-jointed plates swinging chaotically under
gravity while conserving energy. It is visualized in \path{ga_view}
(\path{ga_view/src/active_double_pendulum}, a \path{dynamic_system2dp} advanced via
\path{step(dt)}) and pinned by assertion-checked tests: the planar case in
\path{ga_test/src/ga_appl2dp_physics_test.hpp} and the genuinely-$3D$, non-parallel-axes
case in \path{ga_appl3dp_physics_test.hpp}. Both check energy conservation to RK4
tolerance and the mass-matrix-versus-kinetic-energy identity.}

\paragraph{Spatial pairing and energy:} Each body carries the inertia map $I$ of
subsection~\ref{subsubsection:structure_inertia_maps} (velocity twist $\to$ momentum
wrench $P = I[\Omega]$). A twist is paired with a wrench through the dimension-agnostic
\emph{reciprocal (spatial) product}
\begin{equation}
    \langle \xi, P \rangle = -\,\mathrm{rwdg}(\xi, P),
    \qquad \langle \xi, I[\xi] \rangle = 2\,T
    \label{eq:spatial_pairing}
\end{equation}
i.e. the inertia-map quadratic form equals twice the kinetic energy $T$, carrying the mass
and the angular terms uniformly with no $\tfrac12 m v^2 + \tfrac12 I \omega^2$ split. This
one pairing drives the kinetic energy, the joint-space mass matrix and the forward
dynamics, and lifts to $3D$ by replacing the $3\times3$ map with the $6\times6$ inertia
map and $\mathrm{rwdg}$ with the bivector reciprocal product.

\paragraph{Joints:} A joint is the reduced coordinate linking a body to its parent:
\emph{free} (the unconstrained body), \emph{revolute} (a $1$-DOF hinge) or \emph{prismatic}
(a $1$-DOF slider). Both $1$-DOF joints run identical code. The
body-to-parent motor, the relative twist and the Jacobian column are always
\begin{equation}
    M(q) = M_{\textit{rest}} \veedot \exp_\veedot\!\big(\tfrac12\, q\, S_b\big), \qquad
    \xi = \dot{q}\, S_b, \qquad
    v(\cdot) = \mathrm{rcmt}(S_b, \cdot)
    \label{eq:joint_motor}
\end{equation}
and only the \emph{screw generator} $S_b$ differs: a \emph{finite point} $Q_b$ for a
revolute joint (so $\exp_\veedot$ is a rotation about $Q_b$) versus an \emph{ideal
point}/direction for a prismatic joint (so $\exp_\veedot$ is a translation). The earlier
single-pivot result (eqn.~(\ref{eq:pivot_invariance})) is the special case of one revolute
joint; the generic screw machinery, the prismatic case, and the parallel-axis (Steiner)
inertia emerging automatically from the spatial Jacobian are what the dynamic layer adds.

\paragraph{Forward dynamics:} In reduced coordinates the joint itself \emph{is} the
constraint, so there are no explicit
constraint forces, no drift, and clean energy conservation -- at the cost of a small
linear solve. By virtual work over the bodies one assembles the joint-space equation of
motion $M(q)\,\ddot{q} = \tau(q,\dot{q})$ with
\begin{subeqnarray}
    M_{jk} &=& \sum_i \big\langle\, S_j^{(i)},\; I_i\big[\,S_k^{(i)}\,\big] \,\big\rangle,
               \qquad S_j^{(i)} = \rrev{M_i} \veedot S_j \veedot M_i
               \slabel{eq:fd_mass_matrix} \\
    \tau_j &=& \sum_i m_i \big(\, v_{cm,i}(S_j)\cdot g
               \;-\; v_{cm,i}(S_j)\cdot b_{cm,i} \,\big)
               \slabel{eq:fd_rhs}
\end{subeqnarray}
where the sums run over bodies $i$ that have joint $j$ (respectively $k$) as an ancestor;
$S_j$ is the world joint screw; $S_j^{(i)}$ transports it into body $i$'s frame, where its
inertia map lives; $v_{cm,i}(S_j) = \mathrm{rcmt}(S_j, C_i)$ is the centre-of-mass
velocity per unit joint rate (a column of the spatial Jacobian); $b_{cm,i}$ is the
velocity-product (Coriolis/centripetal) cm-acceleration evaluated at $\ddot{q} = 0$; and
$g$ is gravity. The system is solved for $\ddot{q}$ by a small LU factorization. Each body
still individually obeys the Euler equation~(\ref{eq:ode_system_body_frame}); the mass
matrix is its projection onto the joint coordinates, and the identity
$\tfrac12\,\dot{q}^{\!\top} M(q)\,\dot{q} = T$ ties the two together (a built-in
consistency check). In $3D$ the right-hand side carries the full spatial bias wrench
$I[A_{\textit{bias}}] + \mathrm{rcmt}(V, I[V])$ -- whose second term is the gyroscopic
torque -- which vanishes in the planar case to recover eqn.~(\ref{eq:fd_rhs}).

The system is advanced in time by a Runge-Kutta integrator (RK4): a coupled joint chain is
integrated together in its reduced coordinates, free bodies independently. Energy
conservation to integrator tolerance over the motion -- validated for the free body, the
compound pendulum, the prismatic slider and the chaotic double pendulum, in both $2D$ and
$3D$ -- is the correctness metric used throughout.


\subsubsection{Closed-loop and parallel mechanisms}
\label{subsubsection:closed_loop}

The frame tree of subsection~\ref{subsubsection:frame_trees_systems} reaches every frame
by a \emph{unique} path from the root; that uniqueness is what makes the reduced (joint)
coordinates well-defined. A closed kinematic chain breaks the premise: a four-bar linkage,
a parallelogram, a Stewart-Gough platform each contain a frame reached through several
branches at once, so no single parent assignment describes it. Such loops are admitted
without giving up the energy-clean reduced coordinates by keeping the tree as a
\emph{spanning tree} and re-closing each loop with an explicit constraint between two of
its frames -- \emph{coordinate partitioning}
\cite{Wehage_Haug_coordinate_partitioning:1982}, not the maximal coordinates weighed in
subsection~\ref{subsubsection:addendum_what_ga_buys}. The open-chain assembly is reused
unchanged; the loop closure is an additive layer on top of it
\cite{Featherstone_rigid_body_dynamics:2008}.

\paragraph{Loop-closure constraints:} A loop is closed by requiring two anchor points --
$P_a$ carried by one branch, $P_b$ by another -- to coincide. In $PGA$ coincidence is the
vanishing of the join $P_a \vee P_b$, the line through the two points; realized in
coordinates as the unitized point difference
\begin{equation}
    g(q) = \hat{P}_a(q) - \hat{P}_b(q) = 0,
    \label{eq:loop_residual}
\end{equation}
$m$ scalar equations per loop ($m = 2$ in $2D$, $m = 3$ in $3D$). This is the spatial
spherical (ball) joint that pins a Stewart-Gough leg to its platform, or the pin that
closes a planar four-bar.

\paragraph{Constraint Jacobian:} Differentiating eqn.~(\ref{eq:loop_residual}) along the
motion gives $\dot{g} = G(q)\,\dot{q}$ with $G = \partial g / \partial q$. Each column of
$G$ is the relative partial velocity of the two anchors,
\begin{equation}
    G_{cj} = \frac{\partial (P_a - P_b)}{\partial q_j}
           = \mathrm{rcmt}(S_j, P_a) - \mathrm{rcmt}(S_j, P_b),
    \label{eq:loop_jacobian}
\end{equation}
each term present only when joint $j$ lies on the path to that anchor. But $v(\cdot) =
\mathrm{rcmt}(S_j,\cdot)$ is exactly the spatial-Jacobian column already built for the
mass matrix and the bias force (eqn.~(\ref{eq:fd_mass_matrix}),
eqn.~(\ref{eq:newton_euler_recursion})). The loop closure introduces no new geometry: the
same joint screws $S_j$ that assemble $M$ and $\tau$ assemble $G$.

\paragraph{Constrained equations of motion:} The tree dynamics given by the equation
$M(q)\,\ddot{q} = \tau$ (eqn.~(\ref{eq:fd_mass_matrix})) gain the $m$ constraints $g(q) =
0$. Differentiated twice they read $G\,\ddot{q} = -\dot{G}\,\dot{q}$; adjoining the
constraint forces $G^{\!\top} \lambda$ gives the bordered (KKT) system solved each step,
\begin{equation}
    \begin{bmatrix} M & G^{\!\top} \\[2pt] G & 0 \end{bmatrix}
    \begin{bmatrix} \ddot{q} \\[2pt] \lambda \end{bmatrix}
    =
    \begin{bmatrix} \tau \\[2pt] -\dot{G}\,\dot{q} \end{bmatrix}.
    \label{eq:kkt}
\end{equation}
The Lagrange multipliers $\lambda$ are the constraint wrenches -- the leg forces. $M$ and
$\tau$ are reused unchanged from the open-chain assembly
(eqn.~(\ref{eq:fd_mass_matrix})--(\ref{eq:fd_rhs})); only $G$ and the velocity-product
term $\dot{G}\,\dot{q}$ -- the relative anchor acceleration evaluated at $\ddot{q} = 0$,
read from the same bias pass that yields $b_{cm}$ -- are new. The one constraint,
differentiated zero, once and twice, supplies the three tiers of solve: the position
assembly $g(q) = 0$ by Newton (the dependent joint angles from a driven one), the velocity
distribution $G\,\dot{q} = 0$, and the acceleration distribution $G\,\ddot{q} =
-\dot{G}\,\dot{q}$.

\paragraph{Stabilisation and energy:} Discretely integrated, the position and velocity
constraints drift. They are restored after each step by projection onto the constraint
manifold -- a few Gauss-Newton corrections $q \leftarrow q - G^{+} g$ and $\dot{q}
\leftarrow \dot{q} - G^{+}(G\,\dot{q})$ (the Gear-Gupta-Leimkuhler scheme
\cite{Gear_Gupta_Leimkuhler_constrained_integration:1985}). Because the
constraint forces $G^{\!\top}\lambda$ are orthogonal to the feasible velocities
($G\,\dot{q} = 0$) they do no work, and the projection corrects drift without injecting
energy. The energy-conservation property of the reduced coordinates
(subsection~\ref{subsubsection:addendum_what_ga_buys}) therefore survives the loop closure
-- the objection that ruled out maximal coordinates, whose continuous stabilization trades
energy for feasibility, does not bind on the few constraints a spanning-tree closure adds.
Baumgarte stabilization \cite{Baumgarte_constraint_stabilization:1972} (a term
$2\alpha\,G\dot{q} + \beta^2 g$ on the right-hand side of eqn.~(\ref{eq:kkt})) is the
simpler alternative, kept as a fallback.

\paragraph{The four-bar linkage:} The planar four-bar is the spanning tree crank $\to$
coupler, with the rocker a second branch from the ground, closed by one coincidence of the
coupler tip with the rocker tip: three revolute joints, two equations, one degree of
freedom.\footnote{It is visualized in \path{ga_view}
(\path{ga_view/src/active_four_bar}, a motor-driven \path{closed_loop_system2dp} whose
coupler point traces the classic coupler curve) and pinned by assertion-checked tests: the
planar four-bar in \path{ga_test/src/ga_appl2dp_physics_test.hpp} and a spatial two-arm loop
($3D$) in \path{ga_appl3dp_physics_test.hpp} -- assembly, velocity and acceleration ratios,
energy conservation across a long run, and a bounded closure error $\lVert g\rVert$.}
Driving the
crank, the assembly solves the coupler and rocker angles from
eqn.~(\ref{eq:loop_residual}); the velocity and acceleration ratios follow from
$G\,\dot{q} = 0$ and $G\,\ddot{q} = -\dot{G}\,\dot{q}$ and reproduce the classical
four-bar relations. Released under gravity the single degree of freedom swings, with total
energy conserved to integrator tolerance while $\lVert g\rVert$ stays bounded -- the same
metric the open-chain dynamics use.

\paragraph{The lift to $3D$:} As for the open-chain tier, the step to $3D$ is a change of
representation, not of method. The coincidence residual carries three components instead
of two, the joint screws become lines (\path{twist3dp}, a bivector) and the anchor points
carry the weight $\bv{4}$; eqn.~(\ref{eq:loop_residual})--(\ref{eq:kkt}) and the
projection are unchanged, and the bordered solve is dimension-agnostic. A spatial loop --
two arms meeting at a shared spherical joint, the kinematic core of a Stewart-Gough or
delta platform -- then assembles, distributes its velocities, and conserves energy under
the same machinery (\path{closed_loop_system3dp}).


\subsubsection{The transition from $2D$ to $3D$}
\label{subsubsection:system_2d_to_3d}

Everything above is presented for $2D$ $PGA$ (\path{pga2dp}); the $3D$ stack
(\path{pga3dp}) is the same classes, methods and algorithms, with only the representations
changed:

\begin{table}[H]
\centering\small
\begin{tabular}{|p{0.20\textwidth}|p{0.36\textwidth}|p{0.33\textwidth}|}
    \hline
    \textbf{Concept} & \textbf{\path{pga2dp}} & \textbf{\path{pga3dp}} \\
    \hline \hline
    motor (pose)      & even/odd versor \path{mvec2dp_u} & versor \path{mvec3dp_e} \\
    \hline
    pose orientation  & scalar angle $\varphi$           & axis-angle vector \\
    \hline
    velocity twist    & a vector (grade-$1$)             & a bivector \path{BiVec3dp}
    (grade $2$) \\
    \hline
    point             & vector, weight $\bv{3}=1$        & vector, weight $\bv{4}=1$ \\
    \hline
    joint screw       & a finite/ideal point             & a line / screw axis (a
    bivector) \\
    \hline
    inertia map       & $3\times3$                        & $6\times6$ \\
    \hline
    twist bracket     & $\mathrm{rcmt}(\textit{vec},\textit{vec})$ &
    $\mathrm{rcmt}(\textit{bivec},\textit{bivec})$ \\
    \hline
    spatial pairing   & $-\mathrm{rwdg}(\textit{vec},\textit{bivec})$ &
    $-\mathrm{rwdg}(\textit{bivec},\textit{bivec})$ \\
    \hline
\end{tabular}
\caption{Representation changes in the $2D \to 3D$ lift of the system layer. The motor
         sandwich, the adjoint, the tree walk, the joint exponential, the inertia-map
         pairing and the Euler equation are \emph{identical formulae} in both dimensions.}
\end{table}

The motor sandwich, the adjoint, the tree walk, the joint exponential, the inertia-map
pairing and the Euler equation $\dot{\Omega} = I^{-1}[\,W - \mathrm{rcmt}(\Omega,
I[\Omega])\,]$ are identical formulae in $2D$ and $3D$, differing only in the grade the
twist and wrench occupy. The $2D \to 3D$ step is thus a change of representation, not a
rewrite; this follows from developing the $2D$ case fully first. The same holds for the
closed-loop layer of subsection~\ref{subsubsection:closed_loop}: the loop-closure residual
gains a component and the joint screws become lines, while the constraint Jacobian and the
bordered solve are unchanged.


\subsubsection{Addendum: what geometric algebra contributes to rigid-body dynamics}
\label{subsubsection:addendum_what_ga_buys}

This addendum weighs the approach against classical methods. It is not
part of the modelling exposition; readers interested only in \emph{how} the model works
may skip it. It states where the $PGA$ formulation is advantageous and where a classical
method remains preferable. This library realizes
the \emph{primal} $PGA$ ($\mathbb{G}(2,0,1)$, $\mathbb{G}(3,0,1)$, see
subsection~\ref{subsection:modelling_physics}); the merits below are those of the
geometric formulation as such and hold for the dual plane-based variant too.

The fair baseline is not ``geometric algebra versus nothing'' but geometric algebra versus
the standard tooling: $4\times4$ homogeneous $SE(3)$ matrices, or Featherstone's spatial
vector algebra (six-component twists and wrenches with a hand-defined spatial cross
product) \cite{Featherstone_rigid_body_dynamics:2008}. Both are mature and fast; the
question is whether the reorganization is worthwhile.

\paragraph{Properties of the geometric formulation:}
\begin{enumerate}
    \item \emph{Pose is a motor, not a $4\times4$ matrix.} There is no orthonormality
          constraint to maintain and no renormalization drift; one exponential/logarithm
          covers rotation and translation together as a screw, and the pose evolves on the
          manifold exactly via $M \veedot \exp_\veedot(\tfrac12\,\Omega\,dt)$.
    \item \emph{Twists, momenta and wrenches are bivectors, and the ``spatial cross
          product'' is the algebra's commutator.} A twist is a screw, the
          momentum wrench is the inertia-map image $P = I[\Omega]$, the Newton-Euler bias
          wrench is a bivector, and the spatial cross product is $\mathrm{rcmt}$ -- a
          native product rather than a bespoke operator with conventional signs.
    \item \emph{The adjoint equals the point sandwich.} Transporting a twist between
          frames is the same operation as moving a point
          (eqn.~(\ref{eq:adjoint_is_sandwich})); the classical formulation needs a
          separate $6\times6$ adjoint matrix.
    \item \emph{A force is its line of action -- a wrench.} The combined force-and-moment
          bivector $F = Q \wedge f$
          (subsection~\ref{subsubsection:modelling_force_and_torque}) is exactly the
          object the mechanics literature calls a \emph{wrench}; moments are not a
          separate pseudovector built with an $r \times F$ cross product.
    \item \emph{The Euler equation is one dimension-agnostic line}
          (eqn.~(\ref{eq:ode_system_body_frame})), identical in $2D$ and $3D$ up to grade;
          the $2D \to 3D$ lift is the same formula, not a rewrite.
\end{enumerate}

\paragraph{Coriolis coupling:} The most involved part of a classical double-pendulum
Lagrangian is the
Coriolis/centrifugal coupling -- the Christoffel symbols built from derivatives of the
mass matrix. In the geometric formulation it collapses to one term, the Lie bracket of
twists in the acceleration recursion~(\ref{eq:newton_euler_recursion}): the Coriolis
force is the geometric fact that twists do not commute. These objects are also drawable
in a way six-component vectors are not: each body's instantaneous centre of rotation is a
point read off its world twist, the joint screw is the joint point, and the momentum
wrench is a bivector. The dynamical state is itself geometry, so the
visualization is not a separate construction.

\paragraph{Limitations:} Inertia remains a linear operator ($3\times3$ in $2D$, $6\times6$
in $3D$); for an
anisotropic body it does not collapse to a single multivector. The joint-space mass matrix
$M(q)$ is still a small numerical solve -- reducing to generalized coordinates is a
standard robotics step, and the final $\ddot{q} = M^{-1}\tau$ is ordinary linear algebra.
Performance is not improved either: mature spatial-vector and matrix codes are highly
optimized, so the geometric formulation offers clarity, unification and visualizability
rather than speed.

\paragraph{Reduced vs.\ maximal coordinates:} A more $GA$-native formulation was
deliberately not used: maximal
coordinates, in which every body stays a full motor-and-twist and joints become constraint
wrenches. That formulation is more uniform, but its constraint stabilization
injects and removes energy, which conflicts with an energy-conservation test. The reduced
(joint) coordinates chosen here make the joint itself the constraint, so there are no
explicit constraint forces, no drift, and clean energy -- at the cost of the small $M(q)$
solve. This retains most of the modelling value while conserving energy. The closed-loop
extension (subsection~\ref{subsubsection:closed_loop}) keeps the same choice: a reduced-
coordinate spanning tree plus a few explicit loop constraints, rather than maximal
coordinates throughout, so its constraint forces do no work and the energy behaviour is
preserved.

The recursive Newton-Euler \emph{algorithm} is, of course, classical and not original
here. What geometric algebra contributes is the \emph{representation and the concepts} --
screws, motors, bivectors and the commutator in place of matrices, six-vectors and
hand-defined operators -- together with three properties the standard tooling does not
provide directly: the same primitives scale unchanged from one body to a coupled tree; the
state is visualizable geometry; and the $2D \to 3D$ transition is a change of
representation rather than a rewrite.


\subsubsection{Addendum: further reading}
\label{subsubsection:addendum_further_reading}

The concepts underlying this layer -- motors, the sandwich, the adjoint, twists as a Lie
algebra, the screw unification of rotation and translation, the inertia map, the Coriolis
commutator -- are each standard results with a primary source. The references below
progress from the algebra itself, through the plane-based specialization of $PGA$, to
rigid-body mechanics, and finally to the robotics / Lie-theory and classical viewpoints
that the geometric formulation translates.

A standing caveat applies to every $PGA$ source: this library uses the \emph{primal}
$\mathbb{G}(2,0,1)$/$\mathbb{G}(3,0,1)$ (Lengyel's convention -- points and lines as
primary objects, motions carried by the regressive product), whereas Gunn and Dorst \& De
Keninck use the \emph{dual} plane-based $\mathbb{G}^*$ (planes primary, motions carried by
the geometric product). The two are dual-equivalent
(subsection~\ref{subsection:modelling_physics}); their \emph{concepts} carry over
directly, but their \emph{formulae} are the dual of those used here and must be translated
through the complement relation before being matched against the code.

{\setlength{\emergencystretch}{3em}
\begin{itemize}
    \item \textbf{Geometric algebra foundations.}
    \begin{itemize}
        \item \cite{Macdonald:2010} -- a gentle, rigorous first contact with the geometric
              product, blades and rotors.
        \item \cite{David_Hestenes_NFCM:1986} -- the origin of doing mechanics in
              geometric algebra; its chapter on rigid-body mechanics already points toward
              $PGA$.
        \item \cite{Browne_Grassmann-Algebra_Vol1:2012} -- the exterior-algebra background
              for the wedge and regressive products, complements, and the force and
              momentum lines.
        \item \cite{Perwass_GA_engineering:2009} -- a broad engineering-oriented treatment
              of geometric algebra, covering conformal geometric algebra, numerical and
              uncertainty aspects, and the interactive visualization tool CLUViz.
    \end{itemize}
    \item \textbf{Plane-based $PGA$ and the primal/dual choice.}
    \begin{itemize}
        \item \cite{Lengyel_pga-illuminated:2024} -- the convention this library follows
              (rationale in \cite{Lengyel_poor-foundations_GA:2024} and
              \cite{Lengyel_transwedge_product:2025}).
        \item \cite{Gunn_PGA:2020} and \cite{Plane-based_PGA_Dorst-DeKennik:2022} -- the
              foundational and the most accessible plane-based (dual) treatments; read for
              structure, dualize the formulae.
    \end{itemize}
    \item \textbf{Rigid-body mechanics in $PGA$.}
    \begin{itemize}
        \item \cite{Dynamics_in_PGA_plane-based_Dorst-DeKennik:2023} -- the direct (dual)
              analogue of the dynamics developed here: momentum bivector, the ``forque'',
              the motor equation of motion.
        \item \cite{Gunn_thesis_CayleyKlein:2011} -- the most complete treatment of
              rigid-body kinematics and dynamics with motors and bivector twists in (dual)
              projective geometric algebra.
    \end{itemize}
    \item \textbf{The Lie-theory / robotics translation} (the same groups $SE(2)$/$SE(3)$
          as matrices).
    \begin{itemize}
        \item \cite{Sola_micro_lie_theory:2018} -- the most compact modern account of the
              group/algebra correspondence, the adjoint and the bracket.
        \item \cite{Lynch_Park_modern_robotics:2017} -- the clearest intuition for twists,
              the adjoint and the velocity/acceleration recursion.
        \item \cite{Murray_Li_Sastry_robotic_manipulation:1994} -- the canonical thorough
              treatment.
        \item \cite{Selig_geometric_fundamentals_robotics:2005} -- bridges classical screw
              theory and the Lie-group view.
    \end{itemize}
    \item \textbf{The classical baseline.}
    \begin{itemize}
        \item \cite{Featherstone_rigid_body_dynamics:2008} -- spatial vector algebra:
              six-component twists and wrenches and the spatial cross product, the
              Pl\"ucker-coordinate cousin of the bivector twists and $\mathrm{rcmt}$ used
              here, and the reference for the forward-dynamics algorithm.
    \end{itemize}
    \item \textbf{Constrained multibody dynamics} (the closed-loop layer of
          subsection~\ref{subsubsection:closed_loop}).
    \begin{itemize}
        \item \cite{Wehage_Haug_coordinate_partitioning:1982} -- coordinate partitioning:
              the split into a reduced (spanning-tree) set and the dependent coordinates
              fixed by the loop constraints.
        \item \cite{Gear_Gupta_Leimkuhler_constrained_integration:1985} -- the
              projection (GGL) scheme that restores the position and velocity constraints
              after each step without injecting energy.
        \item \cite{Baumgarte_constraint_stabilization:1972} -- the simpler constraint
              stabilization kept as a fallback.
    \end{itemize}
\end{itemize}
}
\newpage